 \section{Preliminaries}\label{Preliminaries}
 In this section, we recall definitions and basic results for generalized cluster algebras \cite{Chekhov_Shapiro} and generalized quantum cluster algebras \cite{BCDX2018}.

 Throughout this section, we fix a positive integer $n$. Denote by $\mathbb{T}_n$ the $n$-regular tree whose edges are labeled by the numbers $1,\dots,n$ such that the $n$ edges emanating from each vertex carry different labels. We write
 \smash{$
 \xymatrix{t\ar@{-}[r]^k&t'}$}
  to indicate that the vertices $t$, $t'$ of $\mathbb{T}_n$ are linked by an edge labeled by $k$.
For an integer $b$, we use the notation $[b]_+:=\max(b,0)$. We also denote by $[1,n]$ the set $\{1,\dots,n\}$. Denote by $A^{\mathsf T}$ the transpose of a matrix $A$. A non-zero integer vector $\alpha\in \Z^n$ is {\it sign-coherent} if its entries are either all non-negative, or all non-positive.
For an integer vector $\alpha=(a_1,\dots,a_n)^{\mathsf T}\in\mathbb{Z}^n$, we denote by $[\alpha]_+:=([a_1]_+,\dots,[a_n]_+)^{\mathsf T}$. It is clear that $\alpha=[\alpha]_+-[-\alpha]_+$.

\subsection{Generalized cluster algebra}\label{S2}
 We follow \cite{Nakanishi14}.
Let $\mathbb{P}$ be a semifield, whose addition is denoted by $\oplus$. Denote by $\mathbb{ZP}$ the group ring of the multiplicative group $\mathbb{P}$ over $\mathbb{Z}$ and $\mathbb{QP}$ its fraction field. Let
$\mathcal{F}$ be the field of rational function in $n$ variables with coefficients in $\mathbb{QP}$.

\begin{Definition}\label{D4}
 A {\it labeled seed} with coefficients in $\mathbb{P}$ is a triplet $(\mathbf{x},\mathbf{y},B)$ such that
 \begin{itemize}\itemsep=0pt
 \item $B=(b_{ij})_{i,j=1}^{n}$ is a skew-symmetrizable integer matrix;
 \item $\mathbf{x}=(x_1,\dots,x_n)$ is an $n$-tuple of algebraic independent elements of $\mathcal{F}$ over $\mathbb{QP}$;
 \item $\mathbf{y}=(y_1,\dots,y_n)$ is an $n$-tuple of elements in $\mathbb{P}$.

 \end{itemize}
We say that $\mathbf{x}$ is a {\it cluster} and refer to $x_i$, $y_i$ and $B$ as the {\it cluster variables}, the {\it coefficients}
and the {\it exchange matrix}, respectively.
\end{Definition}
For a given labeled seed $(\mathbf{x},\mathbf{y},B)$ and $k\in [1,n]$, we set
\smash{$
\hat{y}_k:=y_k\prod_{j=1}^nx_j^{b_{jk}}
$}
and denote ${\hat{\mathbf{y}}=(\hat{y}_1,\dots, \hat{y}_n)}$. For an integer vector $\mathbf{a}=(a_1,\dots, a_n)^{\mathsf T}\in \Z^n$, we define
\[
\mathbf{x}^\mathbf{a}:=x_1^{a_1}\cdots x_n^{a_n},\qquad \mathbf{y}^\mathbf{a}:=y_1^{a_1}\cdots y_n^{a_n},\qquad \hat{\mathbf{y}}^\mathbf{a}:=\hat{y}_1^{a_1}\cdots \hat{y}_n^{a_n}.
\]

In order to introduce the mutation in generalized cluster algebra, we need the notion of mutation data.
\begin{Definition}\label{D5}
 A {\it mutation data} is a pair $(\mathbf{r},\mathbf{z})$, where
 \begin{itemize}\itemsep=0pt
 \item $\mathbf{r}=(r_1,\dots,r_n)$ is an $n$-tuple of positive integers;
 \item $\mathbf{z}=(z_{i,s})_{i=1,\dots,n;s=1,\dots,r_{i}-1}$ is a family of elements in $\mathbb{P}$ satisfying the reciprocity condition: $z_{i,s}=z_{i,r_i-s}$ for $1\leq s\leq r_{i}-1$.
 \end{itemize}
\end{Definition}
Throughout this subsection, we fix a mutation data $(\mathbf{r},\mathbf{z})$ and set $z_{i,0}=z_{i,r_i}=1$. Now we introduce the $(\mathbf{r},\mathbf{z})$-mutation in generalized cluster algebras.
\begin{Definition}\label{D6}
 For any seed ($\mathbf{x}$,$\mathbf{y}$,$B$) with coefficients in $\mathbb{P}$ and $k\in[1,n]$, the $(\mathbf{r},\mathbf{z})$-mutation of $(\mathbf{x},\mathbf{y},B)$ in direction $k$ is a new seed $\mu_k(\mathbf{x},\mathbf{y},B):=(\mathbf{x}',\mathbf{y}',B') $ with coefficients in $\mathbb{P}$ defined by the following rule:
\begin{align}
x_i'&=
\begin{cases}
 x_i & \text{if $i\neq k$}, \\
 \displaystyle x_k^{-1}\biggl(\prod\limits_{j=1}^{n}x_j^{[-\varepsilon b_{jk}]_+}\biggr)^{r_k}\frac{\sum\limits_{s=0}^{r_k}z_{k,s}\hat{y}_k^{\varepsilon s}}{\mathop{\bigoplus}\limits_{s=0}^{r_k}z_{k,s}y_k^{\varepsilon s}}& \text{if $i=k$},
\end{cases}\label{E4}
\\
 y_i'&=
\begin{cases}
 y_k^{-1} & \text{if $i= k$}, \\
\displaystyle y_i\bigl(y_k^{[\varepsilon b_{ki}]_+}\bigr)^{r_k}\biggl(\mathop{\bigoplus}\limits_{s=0}^{r_k}z_{k,s}y_k^{\varepsilon s}\biggr)^{-b_{ki}}& \text{if $i\neq k$},
\end{cases}\label{E5}
 \\
 b_{ij}'&=
 \begin{cases}
-b_{ij} &\text{if $i=k$ or $j=k$},\\
b_{ij}+r_k([-\varepsilon b_{ik}]_+b_{kj}+b_{ik}[\varepsilon b_{kj}]_+) &\text{else},
\end{cases}\label{E6}
\end{align}
where $\varepsilon\in\{\pm 1\}$.
\end{Definition}

\begin{Remark}\label{R1}\quad
\begin{itemize}\itemsep=0pt
\item[(1)] The mutation formulas \eqref{E4}, \eqref{E5} and \eqref{E6} are independent of the choice of $\varepsilon$ and $\mu_k$ is an involution.
 \item[(2)] If $\mathbf{r}=(1,\dots,1)$, then the mutation formulas \eqref{E4}, \eqref{E5} and \eqref{E6} reduce to the mutation formulas of cluster algebras.
 \item[(3)] The mutation of $\hat{\mathbf{y}}$ is similar as \eqref{E5} of $\mathbf{y}$:
 \begin{align*}
 \hat{y}_i'=
 \begin{cases}
 \hat{y}_k^{-1} & \text{if $i= k$}, \\
 \displaystyle \hat{y}_i\bigl(\hat{y}_k^{[\varepsilon b_{ki}]_+}\bigr)^{r_k}\left(\sum\limits_{s=0}^{r_k}z_{k,s}\hat{y}_k^{\varepsilon s}\right)^{-b_{ki}}& \text{if $i\neq k$}.
 \end{cases}
 \end{align*}
\end{itemize}
\end{Remark}

By assigning the labeled seed $(\mathbf{x},\mathbf{y}, B)$ to a root vertex $t_0\in \mathbb{T}_n$, we obtain an $(\mathbf{r},\mathbf{z})$-seed pattern $t\mapsto \Sigma_t$ of $(\mathbf{x},\mathbf{y}, B)$ in the same way as cluster algebras. In particular, for each vertex~${t\in \mathbb{T}_n}$, we have a labeled seed $\Sigma_t=(\mathbf{x}_t, \mathbf{y}_t, B_t)$ and if \smash{$\xymatrix{t\ar@{-}[r]^k &t'}$}, then $\Sigma_{t'}=\mu_k(\Sigma_t)$. We denote by $\mathbf{x}_t=(x_{1;t},\dots, x_{n;t})$, $\mathbf{y}_t=(y_{1;t},\dots,y_{n;t})$ and $B_t=(\mathbf{b}_{j;t})=(b_{ij;t})$.
\begin{Definition}
 The {\it generalized cluster algebra} $\mathcal{A}:=\mathcal{A}(t\mapsto \Sigma_t)$ associated to the $(\mathbf{r},\mathbf{z})$-seed pattern $t\mapsto \Sigma_t$ of $(\mathbf{x},\mathbf{y}, B)$ is the $\Z \mathbb{P}$-subalgebra of $\mathcal{F}$ generated by
 $\mathcal{X}:=\bigcup_{t\in \mathbb{T}_n}\mathbf{x}_t$.
\end{Definition}
The Laurent phenomenon still holds for generalized cluster algebras.
\begin{Proposition}[{\cite[Theorem 2.5]{Chekhov_Shapiro}}]\label{P6}
 Each cluster variable $x_{i;t}$ could be expressed as a Laurent polynomial of $\mathbf{x}$ with coefficients in $\mathbb{Z}\mathbb{P}$.
\end{Proposition}

\subsection{Generalized cluster algebra with principal coefficients}\label{ss:principal-gca}

From now on, let $\mathbf{y}=(y_1,\dots, y_n)$ and $\mathbf{z}=(z_{i,s})_{i=1,2,\dots,n;s=1,2,\dots,r_i-1}$ with $z_{i,s}=z_{i,r_{i}-s}$ be formal variables and $\mathbb{P}=\operatorname{Trop}(\mathbf{y},\mathbf{z})$ the tropical semifield of $\mathbf{y}$ and $\mathbf{z}$, which is the multiplicative abelian group freely
generated by $\mathbf{y}$ and $\mathbf{z}$ with tropical sum $\oplus$ defined by
\begin{equation*}
\biggl(\prod\limits_{i}y_i^{a_i}\prod\limits_{i,s}z_{i,s}^{a_{i,s}}\biggr)\oplus\biggl(\prod\limits_{i}y_i^{b_i}\prod\limits_{i,s}z_{i,s}^{b_{i,s}}\biggr)=\biggl(\prod\limits_{i}y_i^{\min\{a_i,b_i\}}\prod\limits_{i,s}z_{i,s}^{\min\{a_{i,s},b_{i,s}\}}\biggr),
\end{equation*}
where $a_i,a_{i,s},b_i,b_{i,s}\in \Z$. Let $(\mathbf{x},\mathbf{y}, B)$ be a labeled seed with coefficients in $\mathbb{P}$. Fix an $(\mathbf{r},\mathbf{z})$-seed pattern $t\mapsto \Sigma_t$ of $(\mathbf{x},\mathbf{y}, B)$ by assigning $(\mathbf{x},\mathbf{y}, B)$ to the vertex $t_0\in \mathbb{T}_n$. The associated generalized cluster algebra is called a {\it generalized cluster algebra with principal coefficients}. In this case, we denote it by $\mathcal{A}^\bullet$ to indicate the principal coefficients.


We assign two integer matrices $C_t=(\mathbf{c}_{1;t},\dots,\mathbf{c}_{n;t} )=(c_{ij;t})_{i,j=1}^n$ and $G_t=(\mathbf{g}_{1;t},\dots, \mathbf{g}_{n;t})=(g_{ij;t})_{i,j=1}^n$ to each vertex $t\in \mathbb{T}_n$ by the following recursion:
\begin{itemize}\itemsep=0pt
 \item $C_{t_0}=G_{t_0}=I_n$;
 \item if \smash{$\xymatrix{t\ar@{-}[r]^k&t'}\in \mathbb{T}_n$}, then
 \begin{align}
 c_{ij:t'}&=
 \begin{cases}
 -c_{ij;t} & \text{if $j=k$},\\
 c_{ij;t}+r_k(c_{ik;t}[\varepsilon b_{kj;t}]_++[-\varepsilon c_{ik;t}]_+b_{kj;t})& \text{if $ j\neq k$},
 \end{cases}\label{E9}\\
 \mathbf{g}_{i;t'}&=
 \begin{cases}
 \mathbf{g}_{i;t} & \text{if $i\neq k$},\\
 \displaystyle-\mathbf{g}_{k;t}+r_k\biggl(\sum\limits_{j=1}^{n}[-\varepsilon b_{jk;t}]_+\mathbf{g}_{j;t}-\sum\limits_{j=1}^{n}[-\varepsilon c_{jk;t}]_+\mathbf{b}_{j;t_0}\biggr) & \text{if $i=k$}.
 \end{cases}\label{E10}
 \end{align}
\end{itemize}
We remark that the recurrence formulas \eqref{E9} and \eqref{E10} are independent of the choice of the sign~${\varepsilon\in \{\pm 1\}}$.
We call $t\mapsto C_t$ and $t\mapsto G_t$ the~{\it $(\mathbf{r},\mathbf{z})$-$C$-pattern} and {\it $(\mathbf{r},\mathbf{z})$-$G$-pattern} of the~$(\mathbf{r},\mathbf{z})$-seed pattern of $(\mathbf{x},\mathbf{y}, B)$ respectively.
The column vectors of $C_t$ and $G_t$ are called {\it $c$-vectors} and~{\it $g$-vectors} of the $(\mathbf{r},\mathbf{z})$-seed pattern of $(\mathbf{x},\mathbf{y}, B)$, respectively. We remark that $t\mapsto C_t$ and~${t\mapsto G_t}$ only depend on $B, \mathbf{r}$ and $t_0\in \mathbb{T}_n$.

Denote by $R=\operatorname{diag}\{r_1,\dots,r_n\}$. Note that both $RB$ and $BR$ are skew-symmetrizable. Hence we may assign $(\mathbf{x},\mathbf{y}, RB)$ and $(\mathbf{x},\mathbf{y},BR)$ to the vertex $t_0$ to obtain (ordinary) seed patterns of~$(\mathbf{x},\mathbf{y}, RB)$ and $(\mathbf{x},\mathbf{y},BR)$ respectively.
It was proved by \cite[Proposition 3.9]{Nakanishi14} that the $C$-matrix $C_t$ of the $(\mathbf{r},\mathbf{z})$-seed pattern of $(\mathbf{x},\mathbf{y}, B)$ coincide with the $C$-matrix of the seed pattern of $(\mathbf{x},\mathbf{y}, RB)$ at vertex $t$. Alternatively, $R^{-1}C_tR$ is the $C$-matrix of the ordinary seed pattern of $(\mathbf{x},\mathbf{y}, BR)$ at vertex $t$. As a consequence, every $c$-vector of the $(\mathbf{r},\mathbf{z})$-seed pattern of $(\mathbf{x},\mathbf{y}, B)$ is sign-coherent. In this case, we also say $C_t$ is {\it column sign-coherent}. On the other hand, the $g$-vectors of the $(\mathbf{r},\mathbf{z})$-seed pattern of $(\mathbf{x},\mathbf{y}, B)$ at vertex $t$ coincide with the $g$-vectors of the ordinary seed pattern of $(\mathbf{x},\mathbf{y},BR)$ at vertex $t$. Hence the $G$-matrix $G_t$ is {\it row sign-coherent}, i.e., each row vector of $G_t$ is sign-coherent.
By the sign coherence of $c$-vectors provided in \cite[Corollary 5.5]{GHKK2018}, \eqref{E10} can be rewritten as
\begin{align}\label{E15}
\mathbf{g}_{i;t'}=\begin{cases}\mathbf{g}_{i;t}& \text{if $i\neq k$},\\
 \displaystyle-\mathbf{g}_{k;t}+r_k\biggl(\sum\limits_{j=1}^{n}[-\varepsilon_{k;t} b_{jk;t}]_+\mathbf{g}_{j;t}\biggr)& \text{if $i=k$},
 \end{cases}
\end{align}
where $\varepsilon_{k;t}$ is the common sign of components of the $c$-vector $\mathbf{c}_{k;t}$.

Similar to the ordinary seed pattern, we have the following tropical duality between $C$-matrices and $G$-matrices.
\begin{Proposition}[{\cite[Proposition 3.21]{Nakanishi14}}]\label{P4}
 Let $D_0$ be a diagonal matrix with positive integer diagonal entries such that $D_0RB$ is skew-symmetric. For each $t\in \mathbb{T}_n$, we have
\begin{align}\label{E11}
D_0^{-1}R^{-1}(G_t)^{\mathsf T}D_0RC_t=I_n.
\end{align}
\end{Proposition}
Let $D_0R=\operatorname{diag}\bigl\{d_1^{-1},\dots, d_n^{-1}\bigr\}$. We denote by $(-,-)_{D_0R}\colon \mathbb{Q}^n\times \mathbb{Q}^n\to \mathbb{Q}$ the inner product defined by $(\mathbf{u},\mathbf{v})_{D_0R}=\mathbf{u}^{\mathsf T}D_0R\mathbf{v}$, where $\mathbf{u},\mathbf{v}\in \mathbb{Q}^n$. With this notation, equation \eqref{E11} is equivalent to
\begin{align}\label{E12}
 (\mathbf{g}_{i;t},d_j\mathbf{c}_{j:t})_{D_0R}=\delta_{ij},\qquad \forall i,j\in[1,n].
\end{align}
Furthermore, by noticing that $D_0RB_t$ is skew-symmetric, we have
\begin{align*}
 (\mathbf{u},B_t\mathbf{v})_{D_0R}=-(B_t\mathbf{u},\mathbf{v})_{D_0R},\qquad \forall  \mathbf{u},\mathbf{v}\in \mathbb{Q}^n.
\end{align*}

\begin{Proposition}\label{P5}
 The following equality for the $(\mathbf{r},\mathbf{z})$-seed pattern $t\mapsto \Sigma_t$ holds:
 \begin{align}\label{E14}
 G_tB_t=B_{t_0}C_t.
 \end{align}
\end{Proposition}
\begin{proof}
Since $R^{-1}C_tR$ is the $C$-matrix of the seed pattern of $(\mathbf{x},\mathbf{y}, BR)$ and $G_t$ is the $G$-matrix of the seed pattern of $(\mathbf{x},\mathbf{y}, BR)$, we have
$
G_tB_tR=B_{t_0}R\bigl(R^{-1}C_tR\bigr)
$
by \cite[equation~(6.14)]{fomin_zelevinsky_2007}.
\end{proof}

For the generalized cluster algebra $\mathcal{A}^\bullet$ with principal coefficients, we have the strong Laurent phenomenon.
\begin{Proposition}[{\cite[Proposition 3.3]{Nakanishi14}}]
 Each cluster variable $x_{i;t}$ belongs to $\mathbb{Z}[\mathbf{x}^\pm,\mathbf{y},\mathbf{z}]$.
\end{Proposition}
\begin{Definition}
 The {\it $F$-polynomial} $F_{i;t}:=F_{i;t}[\mathbf{y},\mathbf{z}]$ of the cluster variable $x_{i;t}$ is defined as
 \[
 F_{i;t}[\mathbf{y},\mathbf{z}]:=x_{i;t}|_{x_1=\cdots=x_n=1}\in \mathbb{Z}[\mathbf{y},\mathbf{z}].
 \]
\end{Definition}

\begin{Proposition}[{\cite[Theorem 3.23]{Nakanishi14}}]\label{p:separation-formula-gca}
 For each $t\in \mathbb{T}_n$ and $i\in[1,n]$, the following formula holds:
 $
 x_{i;t}=\mathbf{x}^{\mathbf{g}_{i;t}}F_{i;t}[\hat{\mathbf{y}},\mathbf{z}]$.
\end{Proposition}
\begin{Remark}
 Similar to ordinary cluster algebras, the following statements are equivalent for generalized cluster algebras (cf.\ \cite[Proposition 3.19]{Nakanishi14}):
 \begin{itemize}\itemsep=0pt
 \item every {\it c-vectors} $\mathbf{c}_{i;t}$ is sign-coherent;
 \item every {\it F-polynomial} $F_{i;t}[\mathbf{y},\mathbf{z}]$ has a constant term 1;
 \item every {\it F-polynomial} $F_{i;t}[\mathbf{y},\mathbf{z}]$ has a unique monomial of maximal degree as a polynomial in~$\mathbf{y}$. Moreover its coefficients is 1 and it is divided by all the other occurring monomials as a polynomial in $\mathbf{y}$.
\end{itemize}
\end{Remark}

\subsection{Generalized quantum cluster algebras}\label{ss:gqca}
In this subsection, we introduce the definition and some properties of generalized quantum cluster algebras. We follow \cite{BCDX2018}.

Let $q$ be an indeterminate. Let $m\ge n$ be two positive integers, fix the mutation data $(R,\mathbf{h})$, where $R=\operatorname{diag}\{r_1,\dots,r_n\}$ is a diagonal $n\times n$ matrix whose diagonal coefficients are positive integers and $\mathbf{h}=(\mathbf{h}_1;\dots;\mathbf{h}_n)$ is defined as follows.
For $k\in[1,n]$,
\[
\mathbf{h}_{k}:=\bigl\{h_{k,0}\bigl(q^{\frac{1}{2}}\bigr),h_{k,1}\bigl(q^{\frac{1}{2}}\bigr),\dots,h_{k,r_{k}}\bigl(q^{\frac{1}{2}}\bigr)\bigr\},
\]
where $h_{k,i}\bigl(q^{\frac{1}{2}}\bigr)\in\mathbb{Z}\big[q^{\pm\frac{1}{2}}\big]$ satisfying $h_{k,i}\bigl(q^{\frac{1}{2}}\bigr)=h_{k,r_{k}-i}\bigl(q^{\frac{1}{2}}\bigr)$
and $h_{k,0}\bigl(q^{\frac{1}{2}}\bigr)=h_{k,r_{k}}\bigl(q^{\frac{1}{2}}\bigr)=1$.

A {\it compatible pair} $\bigl(\tilde{B}, \Lambda\bigr)$ consists of an integer $m\times n$-matrix $\tilde{B}$ and a skew-symmetric integer~${m\times m}$-matrix $\Lambda$ such that
\smash{$
\tilde{B}^{\mathsf T}\Lambda=[\begin{matrix} D &  0\end{matrix}]$},
where $D=\operatorname{diag}\bigl\{d_1^{-1},\dots, d_n^{-1}\bigr\}$ is a diagonal~${n\times n}$ matrix whose diagonal coefficients are positive integers. It is easy to see that the principal part~$B$ (i.e., the submatrix formed by the first $n$ rows) of $\tilde{B}$ is skew-symmetrizable and $D$ is a~skew-symmetrizer of $B$.

We define \smash{$E_{k,\varepsilon}^{\tilde{B}R}$} as the $m\times m$-matrix which differs from the identity matrix only in its $k$-th column whose coefficients are given by
 \begin{align*}
 \bigl(E_{k,\varepsilon}^{\tilde{B}R}\bigr)_{ik}=\begin{cases}
 -1& \text{if } i=k,\\
 [-\varepsilon b_{ik}r_k]_+& \text{if } i\neq k.
 \end{cases}
 \end{align*}
 Denote by \smash{$F_{k,\varepsilon}^{R\tilde{B}}$} the $n\times n$-matrix which differs from the identity matrix only in its $k$-th row whose coefficients are given by
 \begin{align*}
 \bigl(F^{R\tilde{B}}_{k,\varepsilon}\bigr)_{ki}=\begin{cases}
 -1& \text{if } i=k,\\
 [\varepsilon r_{k}b_{ki}]_+& \text{if } i\neq k.
 \end{cases}
 \end{align*}
Let $k\in [1,n]$.
The {\it mutation $\mu_k$ in direction $k$} transforms the compatible pair $\bigl(\tilde{B},\Lambda\bigr)$ into \smash{$\mu_{k}\bigl(\tilde{B}, \Lambda\bigr):=\bigl(\tilde{B}', \Lambda'\bigr)$}, where
\[
\tilde{B}'=E_{k,\varepsilon}^{\tilde{B}R}\tilde{B}F_{k,\varepsilon}^{R\tilde{B}},\qquad \Lambda'=\bigl(E_{k,\varepsilon}^{\tilde{B}R}\bigr)^{\mathsf T}\Lambda E_{k,\varepsilon}^{\tilde{B}R}.
\]
It is straightforward to check that the first equality is equivalent to \eqref{E6}. Moreover, $\bigl(\tilde{B}', \Lambda'\bigr)$ is a compatible pair and \smash{$\bigl(\tilde{B}'\bigr)^{\mathsf T}\Lambda'=[D\ 0]$}.

 Fix a skew-symmetric bilinear form $\lambda\colon\mathbb{Z}^{m}\times \mathbb{Z}^m\to \mathbb{Z}$. The {\it quantum torus $\mathcal{T}_\lambda$} associated~with~$\lambda$ is the $\mathbb{Z}\big[q^{\pm\frac{1}{2}}\big]$-algebra generated by the distinguished $\mathbb{Z}\big[q^{\pm\frac{1}{2}}\big]$-basis $\{\mathbf{x}(\alpha)\mid \alpha\in \mathbb{Z}^m\}$ with multiplication given by
\[
\mathbf{x}(\alpha)\mathbf{x}(\beta)=q^{\frac{1}{2}\lambda(\alpha,\beta)}\mathbf{x}(\alpha+\beta)
\]
for any $\alpha,\beta\in \mathbb{Z}^m$. The quantum torus $\mathcal{T}_\lambda$ is an Ore domain and we denote by $\mathcal{F}_q:=\mathcal{F}_\lambda$ its fraction skew field, which will be the ambient field to define the generalized quantum cluster algebras.
\begin{Definition}
 An {\it $(R,\mathbf{h})$-quantum seed} in $\mathcal{F}_q$ is a triple $\Sigma=\bigl(\mathbf{X},\tilde{B},\Lambda\bigr)$, where $\bigl(\tilde{B},\Lambda\bigr)$ is a~compatible pair and $\mathbf{X}=(X_1,\dots, X_m)$ is an $m$-tuple of elements of $\mathcal{F}_q$ such that
 \begin{itemize}\itemsep=0pt
 \item $X_1,\dots, X_{m}$ generated $\mathcal{F}_q$ over $\mathbb{Q}$;
 \item $X_iX_j=q^{\lambda_{ij}}X_jX_i$, where $\Lambda=(\lambda_{ij})$.
 \end{itemize}
 The (labeled) set $\mathbf{X}$ is called a {\it quantum cluster}, $X_1,\dots, X_n$ are {\it quantum cluster variables} and~${X_{n+1},\dots, X_m}$ are coefficients.
\end{Definition}
We define \[\mathbf{X}(\alpha):=q^{\frac{1}{2}\sum_{i<j}a_ia_j\lambda_{ji}}X_1^{a_1}\cdots X_{m}^{a_{m}}\] for any $\alpha=(a_1,\dots, a_{m})^{\mathsf T}\in \mathbb{Z}^{m}$. In particular, $X_i=\mathbf{X}(e_i)$, where $e_1,\dots, e_m$ is the standard $\Z$-basis of $\Z^m$.
It follows that
 \[
 \mathbf{X}(\alpha) \mathbf{X}(\beta)=q^{\frac{1}{2}\alpha^{\mathsf T}\Lambda \beta}\mathbf{X}({\alpha+\beta})
 \]
 for any $\alpha,\beta\in \Z^{m}$. The subalgebra of $\mathcal{F}_q$ generated by $\mathbf{X}(\alpha), \alpha\in \Z^{m}$, is a free $\Z\big[q^{\pm\frac{1}{2}}\big]$-module with basis $\{\mathbf{X}(\alpha) \mid \alpha\in \Z^{m}\}$, and hence identifies with the quantum torus $\mathcal{T}_\Lambda$ associated with the bilinear form $\Lambda\colon\Z^m\times \Z^m\to \Z, (\alpha,\beta)\mapsto \alpha^{\mathsf T}\Lambda \beta$, induced by $\Lambda$.

 For any $\beta\in \Z^n$, we also introduce the notation
$
 \hat{\mathbf{Y}}^{\beta}:=\mathbf{X}\bigl(\tilde{B}\beta\bigr)$.
 It is straightforward to check that
 \[
 \hat{\mathbf{Y}}^{\beta_1}\hat{\mathbf{Y}}^{\beta_2}=q^{\frac{1}{2}\beta_1^{\mathsf T}DB\beta_2}\hat{\mathbf{Y}}^{\beta_1+\beta_2},
 \]
 where $\beta_1,\beta_2\in \Z^n$.

 \begin{Definition}\label{mutation_of_generalzied_seed}
Let $\Sigma=\bigl(\mathbf{X},\tilde{B},\Lambda\bigr)$ be an $(R,\mathbf{h})$-quantum seed in $\mathcal{F}_q$. For any $k\in[1,n]$, the mutation $\mu_k$ in direction $k$ transforms the seed $\Sigma$ into a new triple \smash{$\mu_k(\Sigma):=\bigl(\mathbf{X}',\tilde{B}',\Lambda'\bigr)$}, where
\begin{itemize}\itemsep=0pt
\item[(1)] $\bigl(\tilde{B}', \Lambda'\bigr)=\mu_{k}\bigl(\tilde{B},\Lambda\bigr)$;
\item[(2)] $\mathbf{X}'=(X'_1,\dots, X'_m)$ is given by
\begin{align*}
 X'_i:=\mathbf{X}'(e_i)=\begin{cases}
 \mathbf{X}(e_i)&\text{if }i\neq k,\\
 \displaystyle\sum\limits_{s=0}^{r_k}h_{k,s}\bigl(q^{\frac{1}{2}}\bigr)\mathbf{X}(s[\varepsilon \mathbf{b}_k]_{+}+(r_{k}-s)[-\varepsilon \mathbf{b}_k]_+-e_i)&\text{if }i=k,
 \end{cases}
\end{align*}
where $\mathbf{b}_k$ is the $k$-th column vector of $\tilde{B}$ and $\varepsilon\in\{\pm 1\}$.
\end{itemize}
\end{Definition}
By \cite[Propositions 3.6 and 3.7]{BCDX2018}, the triple $\bigl(\mathbf{X}',\tilde{B}',\Lambda'\bigr)$ is also an $(R,\mathbf{h})$-quantum seed in $\mathcal{F}_q$ and $\mu_{k}$ is an involution.

 For a given $(R,\mathbf{h})$-quantum seed $\bigl(\mathbf{X}, \tilde{B}, \Lambda\bigr)$ in $\mathcal{F}_q$, we assign each vertex $t\in \mathbb{T}_n$ an $(R,\mathbf{h})$-quantum seed $\Sigma_t$ in $\mathcal{F}_q$ which can be obtained from \smash{$\bigl(\mathbf{X},\tilde{B},\Lambda\bigr)$} by iterated mutations such that if~$t$ and $t'$ are linked by an edge labeled $k$, then $\Sigma_{t'}=\mu_k(\Sigma_t)$. We call such an assignment $t\mapsto \Sigma_t$ an {\it $(R,\mathbf{h})$-quantum seed pattern}. It is clear that an $(R,\mathbf{h})$-quantum seed pattern is uniquely determined by the assignment of $\bigl(\mathbf{X},\tilde{B}, \Lambda\bigr)$ to an arbitrary vertex $t_0\in \mathbb{T}_n$. In this case, we refer to $t_0$ the {\it root vertex}, \smash{$\Sigma_{t_0}=\bigl(\mathbf{X}, \tilde{B},\Lambda\bigr)$} the {\it initial $(R,\mathbf{h})$-quantum seed} and $X_1,\dots,X_n$ the {\it initial quantum cluster variables}.
 In the following, when we fix an $(R,\mathbf{h})$-quantum seed pattern, we always denote by $\Sigma_t=\bigl(\mathbf{X}_t,\tilde{B}_t,\Lambda_t\bigr)$ and
\[\mathbf{X}_t=(X_{1;t},\dots,X_{m;t}),\qquad \tilde{B}_t=(b_{ij;t}),\qquad \Lambda_t=(\lambda_{ij;t}).
\]
For the initial quantum seed $\Sigma_{t_0}$, we denote
\[
\mathbf{X}_{t_0}:=\mathbf{X}=(X_1,\dots, X_{m}),\qquad \tilde{B}_{t_0}:=\tilde{B}=(b_{ij}), \qquad\Lambda_{t_0}:=\Lambda=(\lambda_{ij}).
\]
For each vertex $t$, we refer to $\mathbf{X}_t$ a {\it quantum cluster}, $X_{i;t}$ $(1\leq i\leq n)$ {\it quantum cluster variables} and $X_{n+i;t}$ $(1\leq i\leq m-n)$ {\it coefficients}.
It is clear that for every $t\in \mathbb{T}_n$, we have $X_{n+i;t}=X_{n+i}$ for $1\leq i\leq m-n$.
\begin{Definition}
 Given an $(R,\mathbf{h})$-quantum seed pattern $t\mapsto \Sigma_t$, the $(R,\mathbf{h})$-quantum cluster algebra $\mathcal{A}_q(\Sigma_{t_0})$ with initial seed $\Sigma_{t_0}=\bigl(\mathbf{X},\tilde{B},\Lambda\bigr)$ is the \smash{$\mathbb{Z}\big[q^{\pm\frac{1}{2}}\big]\big[X_{n+1}^{\pm 1},\dots,X_m^{\pm 1}\big]$} subalgebra of~$\mathcal{F}_q$ generated by all the quantum cluster variables
 \[\mathcal{X}_q:=\{X_{i;t}\mid 1\le i\le n, t\in \mathbb{T}_n\}.\]
\end{Definition}
\begin{Remark}
The algebra $\mathcal{A}_q(\Sigma_{t_0})$ is also called the {\it generalized quantum cluster algebra} associated with the $(R,\mathbf{h})$-quantum seed $\bigl(\mathbf{X},\tilde{B},\Lambda\bigr)$. In particular, if $R=\operatorname{diag}\{1,\dots, 1\}$, then~$\mathcal{A}_q(\Sigma_{t_0})$ is a quantum cluster algebra in the sense of Berenstein and Zelevinsky \cite{Berenstein_Zelevinsky2005}. On the other hand, it can be viewed as a $q$-deformation for a special generalized cluster algebra in the sense of Chekhov and Shapiro \cite{Chekhov_Shapiro}.
\end{Remark}

Fix an $(R,\mathbf{h})$-quantum seed pattern $t\mapsto \Sigma_t$ with initial seed $\Sigma_{t_0}=\bigl(\mathbf{X},\tilde{B},\Lambda\bigr)$. Denote by~${\mathbf{r}=(r_1,\dots, r_n)}$. Recall that we have the $(\mathbf{r},\mathbf{z})$-$C$-pattern $t\mapsto C_t$ and $(\mathbf{r},\mathbf{z})$-$G$-pattern $t\mapsto G_t$ associated with $B$, $\mathbf{r}$ and $t_0\in \mathbb{T}_n$ as in Section \ref{ss:principal-gca}.

We introduce the $(R,\mathbf{h})$-$G$-pattern $t\mapsto \tilde{G}_t$ for the $(R,\mathbf{h})$-quantum seed pattern $t\mapsto \Sigma_t$ as follows.
 For each vertex $t\in \mathbb{T}_n$, we assign an $m\times m$-integer matrix $\tilde{G}_t=(\tilde{\mathbf{g}}_{1;t},\dots, \tilde{\mathbf{g}}_{m;t})$ to $t$ by the following recursion:
\begin{itemize}\itemsep=0pt
 \item[(1)] $\tilde{G}_{t_0}=I_m$;
 \item[(2)] if \smash{$\xymatrix{t\ar@{-}[r]^k&t'}\in \mathbb{T}_n$}, then
 \begin{align}\label{eq:g-formula}
 \mathbf{\tilde{g}}_{i;t'}=
 \begin{cases}
 \mathbf{\tilde{g}}_{i;t} & \text{if } i\neq k,\\
 \displaystyle -\mathbf{\tilde{g}}_{k;t}+r_{k}\biggl(\sum\limits_{j=1}^{m}[-b_{jk;t}]_+\mathbf{\tilde{g}}_{j;t}-\sum\limits_{j=1}^{n}[- c_{jk;t}]_+\mathbf{b}_{j;t_0}\biggr) & \text{if }i=k.
 \end{cases}
 \end{align}

\end{itemize}
We call $\tilde{G}_t$ the {\it $G$-matrix} of the $(R,\mathbf{h})$-quantum seed pattern $t\mapsto \Sigma_t$ and the column vector $\tilde{\mathbf{g}}_{i;t}$ the {\it $g$-vector} of the quantum cluster variable $X_{i;t}$.
Clearly, we have \smash{$\tilde{G}_t=\left[\begin{smallmatrix}
 G_t&0\\ \star& I_{m-n}
\end{smallmatrix}\right]$}. It follows that
\begin{align}\label{D_form_duality}
 \mathbf{c}^{\mathsf T}_{i;t}[D\quad 0]\mathbf{\tilde{g}}_{j;t}=(\mathbf{c}_{i;t},\mathbf{g}_{j;t})_D=d_{i}^{-1}\delta_{ij}
\end{align}
for $i,j\in[1,n]$.

Similar to \cite[equation (6.14)]{fomin_zelevinsky_2007}, we have the following.
\begin{Proposition}
 For each vertex $t\in \mathbb{T}_n$, the following equation holds:
 \begin{align}\label{duality_quantum_G_C_matrix1}
 \tilde{G}_{t}\tilde{B}_{t}=\tilde{B}_{t_0}C_{t}.
 \end{align}
\end{Proposition}
\begin{proof}

 We prove the equality by induction on the distance between $t_0$ and $t$.
 It is obvious for~${t=t_0}$. Let \smash{$\xymatrix{t\ar@{-}[r]^{k}&t'}$} be an edge in $\mathbb{T}_n$ and suppose that \eqref{duality_quantum_G_C_matrix1} holds for the vertex $t$. We check it for $t'$. We first consider the $k$-th column,
 \begin{align*}
 \begin{aligned}
 \sum\limits^{m}_{i=1}b_{ik;t'}\mathbf{\tilde{g}}_{i;t'}
 &=\sum\limits_{i\neq k}(-b_{ik;t})\mathbf{\tilde{g}}_{i;t}=-\sum\limits^{m}_{i=1}b_{ik;t}\mathbf{\tilde{g}}_{i;t}=-\sum\limits^{n}_{i=1}c_{ik;t}\mathbf{b}_{i;t_0}\qquad
 \text{(by induction)}\\
 &=\sum\limits^{n}_{i=1}c_{ik;t'}\mathbf{b}_{i;t_0}.
 \end{aligned}
 \end{align*}
 Now let $j\neq k$, we have
 \begin{align*}
 \sum\limits^{m}_{i=1}b_{ij;t'}\mathbf{\tilde{g}}_{i;t'}
 ={}&\sum\limits_{i\neq k}b_{ij;t'}\mathbf{\tilde{g}}_{i;t'}+b_{kj;t'}\mathbf{\tilde{g}}_{k;t'}\\
 ={}&\sum\limits_{i\neq k}(b_{ij;t}+r_{k}(b_{ik;t}[b_{kj;t}]_{+}+[-b_{ik;t}]_{+}b_{kj;t}))\mathbf{\tilde{g}}_{i;t}\\
 & +(-b_{kj;t})\left(-\mathbf{\tilde{g}}_{k;t}+r_{k}\left(\sum\limits^{m}_{l=1}[-b_{lk;t}]_{+}\mathbf{\tilde{g}}_{l;t}-\sum\limits^{n}_{l=1}[-c_{lk;t}]_{+}\mathbf{b}_{l;t_0}\right)\right)\\
 ={}&\sum\limits^{m}_{i=1}b_{ij;t}\mathbf{\tilde{g}}_{i;t}+r_{k}[b_{kj;t}]_{+}\sum\limits^{m}_{i=1}b_{ik;t}\mathbf{\tilde{g}}_{i;t}+r_{k}b_{kj;t}\sum\limits^{n}_{i=1}[-c_{ik;t}]_{+}\mathbf{b}_{i;t_0}\\
 ={}&\sum\limits_{i=1}^{n}c_{ij;t}\mathbf{b}_{i;t_0}+r_{k}[b_{kj;t}]_{+}\sum\limits^{n}_{i=1}c_{ik;t}\mathbf{b}_{i;t_0}+r_{k}b_{kj;t}\sum\limits^{n}_{i=1}[-c_{ik;t}]_{+}\mathbf{b}_{i;t_0}\\
 ={}&\sum\limits_{i=1}^{n}(c_{ij;t}+r_{k}(c_{ik;t}[b_{kj;t}]_{+}+[-c_{ik;t}]_{+}b_{kj;t}))\mathbf{b}_{i;t_0}\\
 ={}&\sum\limits_{i=1}^{n}c_{ij;t'}\mathbf{b}_{i;t_0}.
 \end{align*}
 Thus \eqref{duality_quantum_G_C_matrix1} holds for the vertex $t'$. This completes the proof.
\end{proof}

By \eqref{duality_quantum_G_C_matrix1}, formula \eqref{eq:g-formula} is equivalent to the following:
 \begin{align}\label{eq:g-formula-new}
 \mathbf{\tilde{g}}_{i;t'}=
 \begin{cases}
 \mathbf{\tilde{g}}_{i;t} & \text{if } i\neq k,\\
 \displaystyle -\mathbf{\tilde{g}}_{k;t}+r_{k}\biggl(\sum\limits_{j=1}^{m}[-\varepsilon b_{jk;t}]_+\mathbf{\tilde{g}}_{j;t}-\sum\limits_{j=1}^{n}[-\varepsilon c_{jk;t}]_+\mathbf{b}_{j;t_0}\biggr) & \text{if }i=k,
 \end{cases}
 \end{align}
 where $\varepsilon\in\{\pm 1\}$.
By the sign-coherence of $c$-vectors, we have
\begin{align}\label{eq:g-matrix}
 \tilde{G}_{t'}=\tilde{G}_{t}E_{k,\varepsilon_{k;t}}^{\tilde{B}_{t}R},
\end{align}
where $\varepsilon_{k;t}$ is common sign of components of the $c$-vector $\mathbf{c}_{k;t}$.

If $m=2n$ and $\tilde{B}=\left[\begin{smallmatrix}
 B\\ I_n
 \end{smallmatrix}\right]$, then $(R,\mathbf{h})$-quantum cluster algebra $\mathcal{A}_q$ is called a {\it generalized quantum cluster algebra with principal coefficients}. In this case, the $g$-vector $\tilde{\mathbf{g}}_{i;t}$ is closely related to the $g$-vector $\mathbf{g}_{i;t}$.
\begin{Lemma}\label{l:g-vector-comparison}
 Let $m=2n$ and $\tilde{B}=\left[\begin{smallmatrix}
 B\\ I_n
 \end{smallmatrix}\right]$. Then for each vertex $t\in \mathbb{T}_n$ and $1\leq i\leq n$, we have~${\tilde{\mathbf{g}}_{i;t}=\left[\begin{smallmatrix}
 \mathbf{g}_{i;t}\\ 0
 \end{smallmatrix}\right]}$.
\end{Lemma}
\begin{proof}
 Note that $b_{jk;t}=c_{(j-n)k;t}$ whenever $n<j\leq 2n$. By the sign-coherence of $c$-vectors, we may choose a sign $\varepsilon$ such that $[-\varepsilon c_{(n-j)k;t}]_+=0$. Hence the recurrence formula \eqref{eq:g-formula-new} can be rewritten as
 \begin{align*}
 \tilde{\mathbf{g}}_{i;t'}= \begin{cases}
 \mathbf{\tilde{g}}_{i;t} & \text{if } i\neq k,\\
 \displaystyle-\mathbf{\tilde{g}}_{k;t}+r_{k}\sum\limits_{j=1}^{n}[-\varepsilon b_{jk;t}]_+\mathbf{\tilde{g}}_{j;t} & \text{if }i=k.
 \end{cases}
 \end{align*}
 Now the result can be deduced by induction on the distance between $t$ and $t_0$.
\end{proof}

\section[F-polynomials for generalized quantum cluster algebras]{$\boldsymbol{F}$-polynomials for generalized quantum cluster algebras}

Throughout this section, fix an $(R,\mathbf{h})$-quantum seed pattern $t\mapsto \Sigma_t$ with initial seed $\Sigma_{t_0}=\bigl(\mathbf{X},\tilde{B},\Lambda\bigr)$ and keep the notation in Section \ref{ss:gqca}. Recall that $e_1,\dots, e_m$ is the standard $\mathbb{Z}$-basis of $\mathbb{Z}^m$. We denote by $f_1,\dots, f_n$ the standard $\mathbb{Z}$-basis of $\mathbb{Z}^n$.
For an integer vector $\beta=(b_1,\dots, b_m)^{\mathsf T}\in \Z^m$, we denote by $\overline{\beta}=(b_1,\dots, b_n)^{\mathsf T}\in \Z^n$ the truncation of $\beta$.
\subsection{Fock--Goncharov decomposition}\label{s:FG-decom-GQCA}
In this section, we introduce Fock--Goncharov decomposition of mutations for generalized quantum cluster algebras, which generalizes the corresponding construction of quantum cluster algebras in \cite{fock2009,Keller2012}.


For $a,b\in\mathbb{Z}$ and $k\in[1,n]$, let
\begin{align*}
 \left(\sum\limits^{r_{k}}_{s=0}h_{k,s}\bigl(q^{\frac{1}{2}}\bigr)\bigl(q^{\frac{b}{2}}z\bigr)^{s}\right)^{\{a\}}:=
 \begin{cases}\displaystyle\prod\limits_{i=1}^{a}
 \left(\sum\limits^{r_{k}}_{s=0}h_{k,s}\bigl(q^{\frac{1}{2}}\bigr)\bigl(q^{\frac{b(2i-1)}{2}}z\bigr)^{s}\right)&\text{if } a>0,\\
 1&\text{if }a=0,\\
 \displaystyle \prod\limits_{i=a}^{-1}\left(\sum\limits^{r_{k}}_{s=0}h_{k,s}\bigl(q^{\frac{1}{2}}\bigr)\bigl(q^{\frac{b(2i+1)}{2}}z\bigr)^{s}\right)^{-1}&\text{if } a<0.
 \end{cases}
\end{align*}
It is easy to verify that
\begin{gather}
 \left(\sum\limits^{r_{k}}_{s=0}h_{k,s}\bigl(q^{\frac{1}{2}}\bigr)\bigl(q^{\frac{b}{2}}z\bigr)^{s}\right)^{\{a+a'\}}\nonumber\\
 \qquad=\left(\sum\limits^{r_{k}}_{s=0}h_{k,s}\bigl(q^{\frac{1}{2}}\bigr)\bigl(q^{\frac{b}{2}}x\bigr)^{s}\right)^{\{a\}}\bigg|_{x=q^{a'b}z}
 \left(\sum\limits^{r_{k}}_{s=0}h_{k,s}\bigl(q^{\frac{1}{2}}\bigr)\bigl(q^{\frac{b}{2}}z\bigr)^{s}\right)^{\{a'\}}.\label{eq:FG-decom}
\end{gather}

Let $\mathcal{T}_{t}$ be the quantum torus associated with $\Lambda_{t}$, i.e., the $\Z\big[q^{\pm\frac{1}{2}}\big]$-subalgebra of $\mathcal{F}_q$ generated by $\{\mathbf{X}_t(\alpha)\mid \alpha\in \Z^m\}$. It is an Ore domain. Denote by $\mathcal{F}_t$ the fraction skew field of $\mathcal{T}_t$.
Let~\smash{$\xymatrix{t\ar@{-}[r]^k&t'}$} be an edge in $\mathbb{T}_n$. The mutation $\mu_k$ in direction $k$ yields a unique \smash{$\Z\big[q^{\pm\frac{1}{2}}\big]$}-algebra isomorphism $\mu_{k;t}\colon\mathcal{F}_{t'}\to \mathcal{F}_t$ such that
\[
 \mu_{k;t}(\mathbf{X}_{t'}(e_i))=\begin{cases}
 \mathbf{X}_t(e_i) &\text{if $i\neq k$,}\\
 \displaystyle \sum\limits_{s=0}^{r_k}h_{k,s}\bigl(q^{\frac{1}{2}}\bigr)\mathbf{X}_t(s[\varepsilon \mathbf{b}_k]_{+}+(r_{k}-s)[-\varepsilon \mathbf{b}_k]_+-e_i)&\text{if $i=k$.}
 \end{cases}
\]

Recall that for $\alpha\in \Z^n$, we have $\hat{\mathbf{Y}}_t^\alpha=\mathbf{X}_t\bigl(\tilde{B}_t\alpha\bigr)$. For $k\in [1,n]$, we also denote \smash{$\hat{\mathbf{Y}}_{k;t}:=\mathbf{\hat{Y}}_t^{f_k}$}.
\begin{Lemma}
 For $k\in[1,n],t\in\mathbb{T}_{n}$ and $\alpha\in\mathbb{Z}^{n}$, there is a unique $\Z\big[q^{\pm\frac{1}{2}}\big]$-algebra homomorphism~\smash{$\psi_{k;t}\bigl(\mathbf{\hat{Y}}_{t}^{\alpha}\bigr)\colon\mathcal{F}_t\to \mathcal{F}_t$} such that
 \[
\psi_{k;t}\bigl(\mathbf{\hat{Y}}_{t}^{\alpha}\bigr)(\mathbf{X}_t(\beta))=\mathbf{X}_{t}(\beta)\left(\sum\limits_{s=0}^{r_k} h_{k,s}\bigl(q^{\frac{1}{2}}\bigr)\bigl(q^{\frac{1}{2d_{k}}}\mathbf{\hat{Y}}_{t}^{\alpha}\bigr)^{s}\right)^{-\{d_{k}(\bar{\beta},\alpha)_D\}}, \qquad \forall \beta\in \mathbb{Z}^m.
 \]
 \end{Lemma}


 \begin{proof}
It suffices to show that
 \[\psi_{k;t}\bigl(\mathbf{\hat{Y}}_{t}^{\alpha}\bigr)(\mathbf{X}_{t}(\beta_1)\mathbf{X}_{t}(\beta_2))=\psi_{k;t}\bigl(\mathbf{\hat{Y}}_{t}^{\alpha}\bigr)(\mathbf{X}_{t}(\beta_1))\psi_{k;t}\bigl(\mathbf{\hat{Y}}_{t}^{\alpha}\bigr)(\mathbf{X}_{t}(\beta_2))\]
for $\beta_{1},\beta_2\in\mathbb{Z}^{m}$. We have
\begin{align*}
 \psi_{k;t}\bigl(\mathbf{\hat{Y}}_{t}^{\alpha}\bigr)(\mathbf{X}_{t}(\beta_1)\mathbf{X}_{t}(\beta_2))
 ={}&\psi_{k;t}\bigl(\mathbf{\hat{Y}}_{t}^{\alpha}\bigr)\bigl(q^{\frac{1}{2}\beta_{1}^{t}\Lambda_{t}\beta_{2}}\mathbf{X}_{t}(\beta_1+\beta_2)\bigr)\\
 ={}&q^{\frac{1}{2}\beta_{1}^{t}\Lambda_{t}\beta_{2}}\mathbf{X}_{t}(\beta_{1}+\beta_{2})\left(\sum\limits_{s=0}^{r_k} h_{k,s}\bigl(q^{\frac{1}{2}}\bigr)\bigl(q^{\frac{1}{2d_{k}}}\mathbf{\hat{Y}}_{t}^{\alpha}\bigr)^{s}\right)^{-\{d_{k}(\bar{\beta_{1}}+\bar{\beta_{2}},\alpha)_D\}}\\
 ={}&\mathbf{X}_{t}(\beta_{1})\mathbf{X}_{t}(\beta_{2})\left(\sum\limits_{s=0}^{r_k} h_{k,s}\bigl(q^{\frac{1}{2}}\bigr)\bigl(q^{\frac{1}{2d_{k}}}\mathbf{\hat{Y}}_{t}^{\alpha}\bigr)^{s}\right)^{-\{d_{k}(\bar{\beta_{1}},\alpha)+d_k(\bar{\beta_{2}},\alpha)_D\}}\\
={}&\mathbf{X}_{t}(\beta_{1})\mathbf{X}_{t}(\beta_{2})\left(\sum\limits_{s=0}^{r_k} h_{k,s}\bigl(q^{\frac{1}{2}}\bigr)(q^{\frac{1}{2d_{k}}+ (\bar{\beta_{2}},\alpha)_D}\mathbf{\hat{Y}}_{t}^{\alpha})^{s}\right)^{-\{d_{k}(\bar{\beta_{1}},\alpha)_D\}}\\
&\times\left(\sum\limits_{s=0}^{r_k} h_{k,s}\bigl(q^{\frac{1}{2}}\bigr)\bigl(q^{\frac{1}{2d_{k}}}\mathbf{\hat{Y}}_{t}^{\alpha}\bigr)^{s}\right)^{-\{d_{k}(\bar{\beta_{2}},\alpha)_D\}}\\
={}&\mathbf{X}_{t}(\beta_{1})\left(\sum\limits_{s=0}^{r_k} h_{k,s}\bigl(q^{\frac{1}{2}}\bigr)\bigl(q^{\frac{1}{2d_{k}}}\mathbf{\hat{Y}}_{t}^{\alpha}\bigr)^{s}\right)^{-\{d_{k}(\bar{\beta_{1}},\alpha)_D\}}\\
&\times\mathbf{X}_{t}(\beta_{2})\left(\sum\limits_{s=0}^{r_k} h_{k,s}\bigl(q^{\frac{1}{2}}\bigr)\bigl(q^{\frac{1}{2d_{k}}}\mathbf{\hat{Y}}_{t}^{\alpha}\bigr)^{s}\right)^{-\{d_{k}(\bar{\beta_{2}},\alpha)_D\}}\\
={}&\psi_{k;t}\bigl(\mathbf{\hat{Y}}_{t}^{\alpha}\bigr)(\mathbf{X}_{t}(\beta_1))\psi_{k;t}\bigl(\mathbf{\hat{Y}}_{t}^{\alpha}\bigr)(\mathbf{X}_{t}(\beta_2)),
\end{align*}
where the fourth equality follows from equation \eqref{eq:FG-decom}.
\end{proof}

\begin{Lemma}
 For $\alpha\in \Z^n$, $\psi_{k;t}\bigl(\mathbf{\hat{Y}}_{t}^{\alpha}\bigr)$ is an isomorphism. Moreover, its inverse is given by
 \begin{align*}
 \psi_{k;t}\bigl(\mathbf{\hat{Y}}_{t}^{\alpha}\bigr)^{-1}\colon\ \mathcal{F}_t\to \mathcal{F}_t,\qquad
 \mathbf{X}_{t}(\beta)\mapsto\mathbf{X}_{t}(\beta)\left(\sum\limits_{s=0}^{r_k} h_{k,s}\bigl(q^{\frac{1}{2}}\bigr)\bigl(q^{\frac{1}{2d_{k}}}\mathbf{\hat{Y}}_{t}^{\alpha}\bigr)^{s}\right)^{\{d_{k}(\bar{\beta},\alpha)_D\}}.
\end{align*}
\end{Lemma}
\begin{proof}
 It is straightforward to show that there is a unique $\Z\big[q^{\pm\frac{1}{2}}\big]$-algebra homomorphism $\Phi\colon \mathcal{F}_t\to \mathcal{F}_t$ which maps $\mathbf{X}_t(\beta)$ to
 \[
 \mathbf{X}_{t}(\beta)\left(\sum\limits_{s=0}^{r_k} h_{k,s}\bigl(q^{\frac{1}{2}}\bigr)\bigl(q^{\frac{1}{2d_{k}}}\mathbf{\hat{Y}}_{t}^{\alpha}\bigr)^{s}\right)^{\{d_{k}(\bar{\beta},\alpha)_D\}}
 \]
  for any $\beta\in \Z^m$ by \eqref{eq:FG-decom}.
 Furthermore, one can show that
 \[
 \Phi\circ \psi_{k;t}\bigl(\mathbf{\hat{Y}}_{t}^{\alpha}\bigr)(\mathbf{X}_t(\beta))=\mathbf{X}_t(\beta)=\psi_{k;t}\bigl(\mathbf{\hat{Y}}_{t}^{\alpha}\bigr)\circ \Phi(\mathbf{X}_t(\beta))
 \]
 by noticing that $(B_t\alpha, \alpha)_D=0$, which implies the result.
\end{proof}

\begin{Proposition}\label{quantum-decomposition}
 For an edge \smash{$\xymatrix{t\ar@{-}[r]^k&t'}$} in $\mathbb{T}_n$ and $\varepsilon\in \{\pm 1\}$, we have
 \begin{align}\label{eq:FG-dec-gqca}
 \mu_{k;t}=\psi_{k;t}\bigl(\mathbf{\hat{Y}}_{k;t}^{\varepsilon}\bigr)^{\varepsilon}\circ\phi_{k;t;\varepsilon},
 \end{align}
 where $\phi_{k;t;\varepsilon}$ is the unique $\Z\big[q^{\pm\!\frac{1}{2}}\big]$-algebra isomorphism from $\mathcal{F}_{t'}$ to $\mathcal{F}_{\!t}$ taking $\mathbf{X}_{\!t'}(\alpha)$ to $\mathbf{X}_{t}\bigl(E_{\!k,\varepsilon}^{\tilde{B}_{t}R}\alpha\bigr)$ for any $\alpha\in\mathbb{Z}^{m}$.
\end{Proposition}
\begin{proof}
 By definition, we have
 \begin{align*}
 \mathbf{X}_{t'}(e_{k})
 &=\sum\limits^{r_{k}}_{s=0}h_{k;s}\bigl(q^{\frac{1}{2}}\bigr)\mathbf{X}_{t}(s[\varepsilon\mathbf{b}_{k;t}]_{+}+(r_{k}-s)[-\varepsilon\mathbf{b}_{k;t}]_{+}-e_{k})\\
 & =\sum\limits^{r_{k}}_{s=0}h_{k;s}\bigl(q^{\frac{1}{2}}\bigr)\mathbf{X}_{t}(s\varepsilon\mathbf{b}_{k;t}+r_{k}[-\varepsilon\mathbf{b}_{k;t}]_{+}-e_{k})\\
 &=\sum\limits^{r_{k}}_{s=0}h_{k;s}\bigl(q^{\frac{1}{2}}\bigr)\mathbf{X}_{t}\bigl(s\varepsilon\mathbf{b}_{k;t}+E_{k,\varepsilon}^{\tilde{B}_{t}R}e_{k}\bigr)\\
 &=\sum\limits^{r_{k}}_{s=0}h_{k;s}\bigl(q^{\frac{1}{2}}\bigr)q^{\frac{1}{2}\Lambda_{t}(s\varepsilon\mathbf{b}_{k;t},E_{k,\varepsilon}^{\tilde{B}_{t}R}e_{k})}\mathbf{X}_{t}\bigl(E^{\tilde{B}_{t}R}_{k,\varepsilon}e_{k}\bigr)\mathbf{X}_{t}(s\varepsilon\mathbf{b}_{k;t})\\
 &=\sum\limits^{r_{k}}_{s=0}h_{k;s}\bigl(q^{\frac{1}{2}}\bigr)q^{\frac{-s\varepsilon}{2d_{k}}}\mathbf{X}_{t}\bigl(E^{\tilde{B}_{t}R}_{k,\varepsilon}e_{k}\bigr)\mathbf{X}_{t}(s\varepsilon\mathbf{b}_{k;t})\\
 &=\mathbf{X}_{t}\bigl(E^{\tilde{B}_{t}R}_{k,\varepsilon}e_{k}\bigr)\sum\limits^{r_{k}}_{s=0}h_{k;s}\bigl(q^{\frac{1}{2}}\bigr)q^{\frac{-s\varepsilon}{2d_{k}}}\mathbf{X}_{t}(s\varepsilon\mathbf{b}_{k;t})\\
 &=\mathbf{X}_{t}\bigl(E^{\tilde{B}_{t}R}_{k,\varepsilon}e_{k}\bigr)\left(\sum\limits^{r_{k}}_{s=0}h_{k;s}\bigl(q^{\frac{1}{2}}\bigr)\bigl(q^{\frac{1}{2d_{k}}}\mathbf{X}_{t}(\varepsilon\mathbf{b}_{k;t})\bigr)^{s}\right)^{-\varepsilon\{-\varepsilon\}}\\
 &=\mathbf{X}_{t}\bigl(E^{\tilde{B}_{t}R}_{k,\varepsilon}e_{k}\bigr)\left(\sum\limits^{r_{k}}_{s=0}h_{k;s}\bigl(q^{\frac{1}{2}}\bigr)\bigl(q^{\frac{1}{2d_{k}}}\mathbf{\hat{Y}}_{k;t}^{\varepsilon}\bigr)^{s}\right)^{-\varepsilon\{-\varepsilon\}}.
 \end{align*}
 Since \smash{$\bigl(d_{k}\varepsilon f_k,\overline{E^{\tilde{B}_{t}R}_{k,\varepsilon}e_{i}}\bigr)_D=-\varepsilon\delta_{ik}$}, we have
 \begin{align*}
 \mathbf{X}_{t'}(e_{k})=\psi_{k;t}\bigl(\mathbf{\hat{Y}}_{k;t}^{\varepsilon}\bigr)^{\varepsilon}\bigl(\mathbf{X}_{t}\bigl(E^{\tilde{B}_{t}R}_{k,\varepsilon}e_{k}\bigr)\bigr).
 \end{align*}
 For $j\neq k$, we have
 \[
 \mathbf{X}_{t'}(e_{j})
 =\mathbf{X}_{t}(e_{j})
 =\mathbf{X}_{t}\bigl(E^{\tilde{B}_{t}R}_{k,\varepsilon}e_{j}\bigr)
 =\psi_{k;t}\bigl(\mathbf{\hat{Y}}_{k;t}^{\varepsilon}\bigr)^{\varepsilon}\bigl(\mathbf{X}_{t}\bigl(E^{\tilde{B}_{t}R}_{k,\varepsilon}e_{j}\bigr)\bigr).\tag*{\qed}
 \]\renewcommand{\qed}{}
\end{proof}

\begin{Remark}
When $R=I_n$, equation \eqref{eq:FG-dec-gqca} specializes to the Fock--Goncharov decomposition for quantum cluster algebras, where $\phi_{k;t;+}$ and $\phi_{k;t;-}$ are called the {\it tropical parts} of $\mu_{k;t}$ (cf.\ \cite[Section 3.3]{fock2009} and \cite[Section 6.3]{Keller2012}). By setting \smash{$q^{\frac12}=1$}, it further degenerate to the Fock--Goncharov decomposition for cluster algebras, see \cite[Section 2.3]{LMN2023}.
\end{Remark}

Fix a path \smash{$\xymatrix{t_{0}\ar@{-}[r]^{i_1}&t_{1}\ar@{-}[r]^{i_2}&t_{2}\ar@{-}[r]^{i_3}&\cdots\ar@{-}[r]^{i_k}&t_{k}}$} in $\mathbb{T}_n$ and denote it by $\mathbf{i}$. Let $\varepsilon_{j}$ be the common sign of components of $\mathbf{c}_{i_j;t_{j-1}}$ and $\mathbf{c}_{j}^{+}:=\varepsilon_{j}\mathbf{c}_{i_j;t_{j-1}}$ for $j\in[1,k]$.
\begin{Remark}
 The definition of $\varepsilon_j$ is opposite to the one in \cite{LMN2023}, where they define $\varepsilon_j$ to be the common sign of components of $\mathbf{c}_{i_j;t_j}$
 Our convention of $\varepsilon_{j}$ is the same as the one in \cite{Keller2012}, which is more convenient to study maximal green sequences and green-to-red sequences in cluster algebras \cite{KD2020}.
\end{Remark}
Now define
\begin{align*}
&\mu_{t_k}^{t_0}:=\mu_{i_1;t_0}\circ\mu_{i_2;t_1}\circ\cdots\mu_{i_k;t_{k-1}}\colon\ \mathcal{F}_{t_k}\to\mathcal{F}_{t_0},\\
&\psi(\mathbf{i}):=\psi_{i_1;t_0}\bigl(\mathbf{\hat{Y}}_{t_0}^{\mathbf{c}_{1}^+}\bigr)^{\varepsilon_1}\circ\psi_{i_2;t_0}\bigl(\mathbf{\hat{Y}}_{t_0}^{\mathbf{c}_{2}^+}\bigr)^{\varepsilon_2}\circ\cdots\circ\psi_{i_k;t_0}\bigl(\mathbf{\hat{Y}}_{t_0}^{\mathbf{c}_{k}^+}\bigr)^{\varepsilon_k}\colon\ \mathcal{F}_{t_0}\to\mathcal{F}_{t_0}.
\end{align*}
For each $j\in [1,k]$, we also set
\[
\phi_{t_{j}}^{t_0}:=\phi_{i_1;t_{0};\varepsilon_{1}}\circ\phi_{i_2;t_{1};\varepsilon_{2}}\circ\cdots\circ\phi_{i_j;t_{j-1};\varepsilon_{j}}\colon\ \mathcal{F}_{t_{j}}\to\mathcal{F}_{t_{0}}.
\]
It is straightforward to check that $\phi_{t_{j}}^{t_0}$ takes $\mathbf{X}_{t_{j}}(\alpha)$ to $\mathbf{X}_{t_{0}}\bigl(\tilde{G}_{t_{j}}\alpha\bigr)$ by \eqref{eq:g-matrix}, where $\alpha\in \Z^m$.
\begin{Lemma}\label{l:commutative-diag-isom}
 For each $j\in [1,k-1]$, we have
\[
 \phi_{t_{j}}^{t_{0}}\circ\psi_{i_{j+1};t_{j}}\bigl(\mathbf{\hat{Y}}^{\varepsilon_{j+1}}_{i_{j+1};t_{j}}\bigr)^{\varepsilon_{j+1}}
 =\psi_{i_{j+1};t_{0}}\bigl(\mathbf{\hat{Y}}_{t_{0}}^{\mathbf{c}_{j+1}^{+}}\bigr)^{\varepsilon_{j+1}}\circ\phi_{t_{j}}^{t_{0}}.
 \]
\end{Lemma}
\begin{proof}
It is equivalent to show
 \[
 \psi_{i_{j+1};t_{0}}\bigl(\mathbf{\hat{Y}}_{t_{0}}^{\mathbf{c}_{j+1}^{+}}\bigr)\circ\phi_{t_{j}}^{t_{0}}=\begin{cases}
 \phi_{t_{j}}^{t_{0}}\circ\psi_{i_{j+1};t_{j}}\bigl(\mathbf{\hat{Y}}_{i_{j+1};t_{j}}\bigr) & \text{if $\varepsilon_{j+1}=+1$},\\
 \phi_{t_{j}}^{t_{0}}\circ\psi_{i_{j+1};t_{j}}\bigl(\mathbf{\hat{Y}}^{-1}_{i_{j+1};t_{j}}\bigr) & \text{if $\varepsilon_{j+1}=-1$.}
 \end{cases}
 \]
 For any $\alpha=(\alpha_{1},\dots,\alpha_{m})^{\mathsf T}\in\mathbb{Z}^{m}$,
 \begin{gather*}
  \phi_{t_{j}}^{t_{0}}\circ\psi_{i_{j+1};t_{j}}\bigl(\mathbf{\hat{Y}}_{i_{j+1};t_{j}}^{\varepsilon_{j+1}}\bigr)(\mathbf{X}_{t_{j}}(\alpha))\\
  \qquad=\phi_{t_{j}}^{t_{0}}\left(\mathbf{X}_{t_{j}}(\alpha)\left(\sum\limits_{s=0}^{r_{i_{j+1}}} h_{i_{j+1},s}\bigl(q^{\frac{1}{2}}\bigr)\bigl(q^{\frac{1}{2d_{i_{j+1}}}}\mathbf{\hat{Y}}_{i_{j+1};t_{j}}^{\varepsilon_{j+1}}\bigr)^{s}\right)^{-\{d_{i_{j+1}}(\bar{\alpha},\varepsilon_{j+1}f_{i_{j+1}})_D\}}\right)\\
 \qquad =\mathbf{X}_{t_{0}}\bigl(\tilde{G}_{t_{j}}\alpha\bigr)\left(\sum\limits_{s=0}^{r_{i_{j+1}}} h_{i_{j+1},s}\bigl(q^{\frac{1}{2}}\bigr)\bigl(q^{\frac{1}{2d_{i_{j+1}}}}\mathbf{X}_{t_0}\bigl(\varepsilon_{j+1}\tilde{G}_{t_{j}}\tilde{B}_{t_{j}}f_{i_{j+1}}\bigr)\bigr)^{s}\right)^{-\{d_{i_{j+1}}(\bar{\alpha},\varepsilon_{j+1}f_{i_{j+1}})_D\}}\\
 \qquad =\mathbf{X}_{t_{0}}\bigl(\tilde{G}_{t_{j}}\alpha\bigr)\left(\sum\limits_{s=0}^{r_{i_{j+1}}} h_{i_{j+1},s}\bigl(q^{\frac{1}{2}}\bigr)\bigl(q^{\frac{1}{2d_{i_{j+1}}}}\mathbf{X}_{t_0}\bigl(\varepsilon_{j+1}\tilde{B}_{t_{0}}C_{t_{j}}f_{i_{j+1}}\bigr)\bigr)^{s}\right)^{-\{d_{i_{j+1}}(\bar{\alpha},\varepsilon_{j+1}f_{i_{j+1}})_D\}}\\
 \qquad =\mathbf{X}_{t_{0}}\bigl(\tilde{G}_{t_{j}}\alpha\bigr)\left(\sum\limits_{s=0}^{r_{i_{j+1}}} h_{i_{j+1},s}\bigl(q^{\frac{1}{2}}\bigr)\bigl(q^{\frac{1}{2d_{i_{j+1}}}}\mathbf{X}_{t_0}\bigl(\tilde{B}_{t_{0}}\mathbf{c}_{j+1}^{+}\bigr)\bigr)^{s}\right)^{-\{d_{i_{j+1}}(\bar{\alpha},\varepsilon_{j+1}f_{i_{j+1}})_D\}}\\
 \qquad =\mathbf{X}_{t_{0}}\bigl(\tilde{G}_{t_{j}}\alpha\bigr)\left(\sum\limits_{s=0}^{r_{i_{j+1}}} h_{i_{j+1},s}\bigl(q^{\frac{1}{2}}\bigr)\bigl(q^{\frac{1}{2d_{i_{j+1}}}}\mathbf{\hat{Y}}_{t_{0}}^{\mathbf{c}_{j+1}^{+}}\bigr)^{s}\right)^{-\{\alpha_{i_{j+1}}\varepsilon_{j+1}\}},
 \end{gather*}
 where the third equality follows from \eqref{duality_quantum_G_C_matrix1}.
 On the other hand,
 \begin{gather*}
 \psi_{i_{j+1};t_{0}}\bigl(\mathbf{\hat{Y}}_{t_{0}}^{\mathbf{c}_{j+1}^{+}}\bigr)\circ\phi_{t_{j}}^{t_{0}}(\mathbf{X}_{t_{j}}(\alpha))\\
\qquad =\psi_{i_{j+1};t_{0}}\bigl(\mathbf{\hat{Y}}_{t_{0}}^{\mathbf{c}_{j+1}^{+}}\bigr)\bigl(\mathbf{X}_{t_{0}}\bigl(\tilde{G}_{t_{j}}\alpha\bigr)\bigr)\\
\qquad =\mathbf{X}_{t_{0}}\bigl(\tilde{G}_{t_{j}}\alpha\bigr)\left(\sum\limits_{s=0}^{r_{i_{j+1}}} h_{i_{j+1},s}\bigl(q^{\frac{1}{2}}\bigr)\bigl(q^{\frac{1}{2d_{i_{j+1}}}}\mathbf{\hat{Y}}_{t_{0}}^{\mathbf{c}_{j+1}^{+}}\bigr)^{s}\right)^{-\{d_{i_{j+1}}(\overline{\tilde{G}_{t_{j}}\alpha},\mathbf{c}_{j+1}^{+})_D\}}\\
\qquad =\mathbf{X}_{t_{0}}\bigl(\tilde{G}_{t_{j}}\alpha\bigr)\left(\sum\limits_{s=0}^{r_{i_{j+1}}} h_{i_{j+1},s}\bigl(q^{\frac{1}{2}}\bigr)\bigl(q^{\frac{1}{2d_{i_{j+1}}}}\mathbf{\hat{Y}}_{t_{0}}^{\mathbf{c}_{j+1}^{+}}\bigr)^{s}\right)^{-\{\alpha_{i_{j+1}}\varepsilon_{j+1}\}}\qquad\text{(by \eqref{D_form_duality})}.
 \end{gather*}
 This completes the proof.
\end{proof}

\begin{Proposition}\label{long-quantum-decomposition}
 Keep the notation as above, we have
 $\mu_{t_k}^{t_0}=\psi(\mathbf{i})\circ\phi_{t_{k}}^{t_0}\colon  \mathcal{F}_{t_{k}}\to\mathcal{F}_{t_{0}}$.
\end{Proposition}
\begin{proof}
 By using Proposition \ref{quantum-decomposition} and Lemma \ref{l:commutative-diag-isom}, we have
 \begin{align*}
 \mu_{t_k}^{t_0}
 ={}&\mu_{i_1;t_0}\circ\mu_{i_2;t_1}\circ\cdots\circ\mu_{i_k;t_{k-1}}\\
 ={}&\psi_{i_{1};t_{0}}\bigl(\mathbf{\hat{Y}}_{i_{1};t_{0}}^{\varepsilon_{1}}\bigr)^{\varepsilon_{1}}\circ\phi_{i_{1};t_{0};\varepsilon_{1}}\circ\psi_{i_{2};t_{1}}
 \bigl(\mathbf{\hat{Y}}_{i_{2};t_{1}}^{\varepsilon_{2}}\bigr)^{\varepsilon_{2}}\circ\phi_{i_{2};t_{1};\varepsilon_{2}}\\
 &\circ\cdots\circ\psi_{i_{k};t_{k-1}}
 \bigl(\mathbf{\hat{Y}}_{i_{k};t_{k-1}}^{\varepsilon_{k}}\bigr)^{\varepsilon_{k}}\circ\phi_{i_{k};t_{k-1};\varepsilon_{k}}\\
 ={}&\psi_{i_{1};t_{0}}\bigl(\mathbf{\hat{Y}}_{t_{0}}^{\mathbf{c}_{1}^{+}}\bigr)^{\varepsilon_{1}}\circ\psi_{i_{2};t_{1}}\bigl(\mathbf{\hat{Y}}_{t_{0}}^{\mathbf{c}_{2}^{+}}\bigr)^{\varepsilon_{2}}\circ\phi_{i_{1};t_{0};\varepsilon_{1}}\circ\phi_{i_{2};t_{1};\varepsilon_{2}}\\
 &\circ\cdots\circ\psi_{i_{k};t_{k-1}}(\mathbf{\hat{Y}}_{i_{k};t_{k-1}}^{\varepsilon_{k}})^{\varepsilon_{k}}\circ\phi_{i_{k};t_{k-1};\varepsilon_{k}}\\
 &\vdots\\
 ={}&\psi_{i_{1};t_{0}}\bigl(\mathbf{\hat{Y}}_{t_{0}}^{\mathbf{c}_{1}^{+}}\bigr)^{\varepsilon_{1}}\circ\psi_{i_{2};t_{1}}\bigl(\mathbf{\hat{Y}}_{t_{0}}^{\mathbf{c}_{2}^{+}}\bigr)^{\varepsilon_{2}}\circ\cdots\circ\psi_{i_{k};t_{0}}
 \bigl(\mathbf{\hat{Y}}_{t_{0}}^{\mathbf{c}_{k}^{+}}\bigr)^{\varepsilon_{k}}\circ\phi_{t_{k}}^{t_0}\\
 ={}&\psi(\mathbf{i})\circ\phi_{t_{k}}^{t_0}.\tag*{\qed}
 \end{align*}\renewcommand{\qed}{}
\end{proof}

\subsection[F-polynomials]{$\boldsymbol{F}$-polynomials}\label{s:F-poly_gen_quantum_cluster_alg}

 In this section, we define $F$-polynomials for the generalized quantum cluster algebra $\mathcal{A}_q(\Sigma_{t_0})$ and prove their polynomial property under a mild condition.
Recall that for $i,j\in[1,n]$, we have
\begin{align}\label{eq:quasi-com}
\mathbf{\hat{Y}}_{i,t_{0}}\mathbf{\hat{Y}}_{j,t_{0}}=q^{d_{i}^{-1}b_{ij}}\mathbf{\hat{Y}}_{j,t_{0}}\mathbf{\hat{Y}}_{i,t_{0}}.
\end{align}
The quasi-commutative relations \eqref{eq:quasi-com} depend only on the entries of $D$ and the principal part~$B$ of $\tilde{B}_{t_0}$.
Let $\mathcal{T}_{DB}$ be the quantum torus associated to $DB$ whose underlying space is the free~\smash{$\mathbb{Z}\big[q^{\frac{1}{2}}\big]$}-module with basis $\{\mathbf{Z}(\alpha)| \alpha\in\mathbb{Z}^{n}\}$, and its multiplication is given by
\[\mathbf{Z}(\alpha)\mathbf{Z}(\beta)=q^{\frac{1}{2}\alpha^{\mathsf T}DB\beta}\mathbf{Z}(\alpha+\beta).\]
Denote by $\mathcal{F}_{DB}$ the fraction skew field of $\mathcal{T}_{DB}$ and $\mathbf{Z}_{i}=\mathbf{Z}(f_{i})$.

Recall that \smash{$\xymatrix{t_{0}\ar@{-}[r]^{i_1}&t_{1}\ar@{-}[r]^{i_2}&t_{2}\ar@{-}[r]^{i_3}&\cdots\ar@{-}[r]^{\!\!i_k}&t_{k}}$} is a path in $\mathbb{T}_n$, and $\varepsilon_{j}$ is the common sign of components of $\mathbf{c}_{i_j;t_{j-1}}$. For simplicity of notation, we also denote by
\begin{alignat*}{6}
& d_{(j)}=d_{i_j},\qquad&&
r_{(j)}=r_{i_j},  \qquad &&  \mathbf{c}_{j}=\mathbf{c}_{i_j;t_{j-1}},\qquad&&    \mathbf{c}_{j}^+=\varepsilon_{j}\mathbf{c}_{j}, \qquad&&   \mathbf{\hat{c}}_{j}^+=B\mathbf{c}_{j}^+,& \\
& \mathbf{\tilde{g}}_{j}=\mathbf{\tilde{g}}_{i_j;t_j}, \qquad&&   \mathbf{g}_{j}=\overline{\mathbf{\tilde{g}}_{j}}.&&&&&&&
\end{alignat*}

We first define a set of elements $\{ L_{i,j}\mid i,j\in [1,k]\}$ of $\mathcal{F}_q$ by the initial condition
\[L_{1,i}:=\mathbf{\hat{Y}}_{t_{0}}^{\mathbf{c}_{i}^{+}}\left(\sum\limits^{r_{(1)}}_{s=0}h_{i_{1},s}
\bigl(q^{\frac{1}{2}}\bigr)\bigl(q^{\frac{1}{2d_{(1)}}}\mathbf{\hat{Y}}_{t_0}^{\mathbf{c}_{1}^{+}}\bigr)^{s}\right)^{-\varepsilon_{1}\{(d_{(1)}\mathbf{c}_{1}^{+},\mathbf{\hat{c}}_{i}^{+})_{D}\}} \qquad \text{for}\  i\in [1,k]
\]
with recurrence relations
\[
 L_{j+1,i}=L_{j,i}\left(\sum\limits^{r_{(j+1)}}_{s=0}h_{i_{j+1},s}\bigl(q^{\frac{1}{2}}\bigr)
 \bigl(q^{\frac{1}{2d_{(j+1)}}}L_{j,j+1}\bigr)^{s} \right)^{-\varepsilon_{j+1}\{(d_{(j+1)}\mathbf{c}_{j+1}^{+},\mathbf{\hat{c}}_{i}^{+})_{D}\}} \quad \text{for}\  j \in [1,k-1].
\]
Then set
\begin{align*}
 &L_1=\sum\limits^{r_{(1)}}_{s=0}h_{i_{1},s}\bigl(q^{\frac{1}{2}}\bigr)\bigl(q^{\frac{1}{2d_{(1)}}}\mathbf{\hat{Y}}_{t_{0}}^{\mathbf{c}_{1}^{+}}\bigr)^{s},\\
 &L_{l+1}=\sum\limits_{s=0}^{r_{(l+1)}}h_{i_{l+1},s}\bigl(q^{\frac{1}{2}}\bigr)\bigl(q^{\frac{1}{2d_{(l+1)}}}L_{l,l+1}\bigr)^{s},\qquad l\in[1,k-1].
 \end{align*}


\begin{Lemma}\label{product_of_F_polynomial}
 Keep the notation as above. We have
 \begin{align}\label{eq:f-poly-well-def}
 \mathbf{X}_{t_{0}}(-\mathbf{\tilde{g}}_{k})\mu_{t_{k}}^{t_{0}}(\mathbf{X}_{i_{k};t_{k}})=\prod\limits^{\longrightarrow}_{j\in[1,k]}L_{j}^{-\varepsilon_{j}\{d_{(j)}(\mathbf{c}_{j}^{+},\mathbf{g}_{k})_D\}}.
 \end{align}

\end{Lemma}
\begin{proof}
 Let $\mathbf{i}_{j}$ be the sub-sequence \smash{$\xymatrix{t_{0}\ar@{-}[r]^{i_1}&t_{1}\ar@{-}[r]^{i_2}&t_{2}\ar@{-}[r]^{i_3}&\cdots\ar@{-}[r]^{i_j}&t_{j}}$} for $j\in [1,k]$.
 By Proposition \ref{long-quantum-decomposition},
 \begin{align*}
 \mathbf{X}_{t_{0}}(-\mathbf{\tilde{g}}_{k})\mu_{t_{k}}^{t_{0}}(\mathbf{X}_{i_{k};t_{k}})
 ={}&\mathbf{X}_{t_{0}}(-\mathbf{\tilde{g}}_{k})\psi(\mathbf{i})\circ\phi_{t_{k}}^{t_0}(\mathbf{X}_{i_{k};t_{k}})\\
 ={}&\mathbf{X}_{t_{0}}(-\mathbf{\tilde{g}}_{k})\psi(\mathbf{i})(\mathbf{X}_{t_{0}}(\mathbf{\tilde{g}}_{k}))\\
 ={}&\mathbf{X}_{t_{0}}(-\mathbf{\tilde{g}}_{k})\psi(\mathbf{i}_{1})(\mathbf{X}_{t_{0}}(\mathbf{\tilde{g}}_{k}))\psi(\mathbf{i}_{1})(\mathbf{X}_{t_{0}}(-\mathbf{\tilde{g}}_{k}))\psi(\mathbf{i}_{2})(\mathbf{X}_{t_{0}}(\mathbf{\tilde{g}}_{k}))\\
 &\cdots\psi(\mathbf{i}_{k-1})(\mathbf{X}_{t_{0}}(-\mathbf{\tilde{g}}_{k}))\psi(\mathbf{i}_{k})(\mathbf{X}_{t_{0}}(\mathbf{\tilde{g}}_{k})).
 \end{align*}
 We first claim that $L_{j,i}=\psi(\mathbf{i}_{j})\bigl(\mathbf{\hat{Y}}_{t_{0}}^{\mathbf{c}_{i}^{+}}\bigr)$
 for $i,j\in[1,k]$. We prove it by induction on $j$.
 For $j=1$, we have
 \begin{align*}
 \psi(\mathbf{i}_{1})\bigl(\mathbf{\hat{Y}}_{t_{0}}^{\mathbf{c}_{i}^{+}}\bigr)
 &=\psi_{i_{1}:t_{0}}\bigl(\mathbf{\hat{Y}}_{t_{0}}^{\mathbf{c}_{1}^{+}}\bigr)^{\varepsilon_{1}}\bigl(\mathbf{\hat{Y}}_{t_{0}}^{\mathbf{c}_{i}^{+}}\bigr)=\mathbf{\hat{Y}}_{t_{0}}^{\mathbf{c}_{i}^{+}}\left(\sum\limits^{r_{(1)}}_{s=0}h_{i_{1},s}\bigl(q^{\frac{1}{2}}\bigr)
 \bigl(q^{\frac{1}{2d_{(1)}}}\mathbf{\hat{Y}}_{t_{0}}^{\mathbf{c}_{1}^{+}}\bigr)^{s}\right)^{-\varepsilon_{1}\{(\mathbf{\hat{c}}_{i}^{+},d_{(1)}\mathbf{c}_{1}^{+})_{D}\}}\\
 &=L_{1,i}.
 \end{align*}
 Suppose that $L_{j,i}=\psi(\mathbf{i}_{j})\bigl(\mathbf{\hat{Y}}_{t_{0}}^{\mathbf{c}_{i}^{+}}\bigr)$ for any $i\in[1,k]$, then
 \begin{align*}
 \psi(\mathbf{i}_{j+1})\bigl(\mathbf{\hat{Y}}_{t_{0}}^{\mathbf{c}_{i}^{+}}\bigr)
 ={}&\psi(\mathbf{i}_{j})\bigl(\psi_{i_{j+1};t_{0}}\bigl(\mathbf{\hat{Y}}_{t_{0}}^{\mathbf{c}_{j+1}^{+}}\bigr)^{\varepsilon_{j+1}}\bigl(\mathbf{\hat{Y}}_{t_{0}}^{\mathbf{c}_{i}^{+}}\bigr)\bigr)\\
 ={}&\psi(\mathbf{i}_{j})\left(\bigl(\mathbf{\hat{Y}}_{t_{0}}^{\mathbf{c}_{i}^{+}}\bigr)\left(\sum\limits^{r_{(j+1)}}_{s=0}h_{i_{j+1},s}
 \bigl(q^{\frac{1}{2}}\bigr)\bigl(q^{\frac{1}{2d_{(j+1)}}}\mathbf{\hat{Y}}_{t_{0}}^{\mathbf{c}_{j+1}^{+}}\bigr)^{s}\right)^{-\varepsilon_{j+1}
 \{(\mathbf{\hat{c}}_{i}^{+},d_{(j+1)}\mathbf{c}_{j+1}^{+})_{D}\}}\right)\\
 ={}&\psi(\mathbf{i}_{j})\bigl(\mathbf{\hat{Y}}_{t_{0}}^{\mathbf{c}_{i}^{+}}\bigr)\\
 &\times\left(\sum\limits^{r_{(j+1)}}_{s=0}h_{i_{j+1},s}\bigl(q^{\frac{1}{2}}\bigr)
 \bigl(q^{\frac{1}{2d_{(j+1)}}}\psi(\mathbf{i}_{j})\bigl(\mathbf{\hat{Y}}_{t_{0}}^{\mathbf{c}_{j+1}^{+}}\bigr)\bigr)^{s}\right)^{-\varepsilon_{j+1}\{(\mathbf{\hat{c}}_{i}^{+},d_{(j+1)}\mathbf{c}_{j+1}^{+})_{D}\}}\\
 ={}&L_{j,i}\left(\sum\limits^{r_{(j+1)}}_{s=0}h_{i_{j+1},s}\bigl(q^{\frac{1}{2}}\bigr)\bigl(q^{\frac{1}{2d_{(j+1)}}}L_{j,j+1}\bigr)^{s}\right)^{-\varepsilon_{j+1}
 \{(\mathbf{\hat{c}}_{i}^{+},d_{(j+1)}\mathbf{c}_{j+1}^{+})_{D}\}}\\
 ={}&L_{j+1,i}.
 \end{align*}
 This completes the proof of the claim.
 A direct computation shows that
 \begin{align*}
 \mathbf{X}_{t_0}(-\mathbf{\tilde{g}}_{k})\psi(\mathbf{i}_{1})(\mathbf{X}_{t_0}(\mathbf{\tilde{g}}_{k}))
 &=\mathbf{X}_{t_0}(-\mathbf{\tilde{g}}_{k})\psi_{i_{1};t_{0}}\bigl(\mathbf{\hat{Y}}_{t_{0}}^{\mathbf{c}_{1}^{+}}\bigr)^{\varepsilon_{1}}(\mathbf{X}_{t_0}(\mathbf{\tilde{g}}_{k}))\\
 &=\left(\sum\limits^{r_{(1)}}_{s=0}h_{i_{1},s}\bigl(q^{\frac{1}{2}}\bigr)\bigl(q^{\frac{1}{2d_{(1)}}}\mathbf{\hat{Y}}_{t_{0}}^{\mathbf{c}_{1}^{+}}\bigr)^{s}\right)^{-\varepsilon_{1}\{(\mathbf{g}_{k},d_{(1)}\mathbf{c}_{1}^{+})_{D}\}}\\
 &=L_{1}^{-\varepsilon_{1}\{(\mathbf{g}_{k},d_{(1)}\mathbf{c}_{1}^{+})_{D}\}}.
 \end{align*}
 For $j\in[1,k-1]$, we have
 \begin{gather*}
  \psi(\mathbf{i}_{j})(\mathbf{X}_{t_{0}}(-\mathbf{\tilde{g}}_{k}))\psi(\mathbf{i}_{j+1})(\mathbf{X}_{t_{0}}(\mathbf{\tilde{g}}_{k}))\\
  \qquad=\psi(\mathbf{i}_{j})\bigl(\mathbf{X}_{t_{0}}(-\mathbf{\tilde{g}}_{k})\psi_{i_{j+1};t_{j}}\bigl(\mathbf{\hat{Y}}_{t_0}^{\mathbf{c}_{j+1}^+}\bigr)^{\varepsilon_{j+1}}
  (\mathbf{X}_{t_{0}}(\mathbf{\tilde{g}}_{k}))\bigr)\\
 \qquad  =\psi(\mathbf{i}_{j})\left(\left(\sum\limits^{r_{(j+1)}}_{s=0}h_{i_{j+1},s}\bigl(q^{\frac{1}{2}}\bigr)
 \bigl(q^{\frac{1}{2d_{(j+1)}}}\mathbf{\hat{Y}}_{t_{0}}^{\mathbf{c}_{j+1}^{+}}\bigr)^{s}\right)
 ^{-\varepsilon_{j+1}\{(\mathbf{g}_{k},d_{(j+1)}\mathbf{c}_{j+1}^{+})_{D}\}}\right)\\
\qquad  =\left(\sum\limits^{r_{(j+1)}}_{s=0}h_{i_{j+1},s}\bigl(q^{\frac{1}{2}}\bigr)\bigl(q^{\frac{1}{2d_{(j+1)}}}\psi(\mathbf{i}_{j})\bigl(\mathbf{\hat{Y}}_{t_{0}}^{\mathbf{c}_{j+1}^{+}}\bigr)\bigr)^{s}\right)^{-\varepsilon_{j+1}\{(\mathbf{g}_{k},d_{(j+1)}\mathbf{c}_{j+1}^{+})_{D}\}}\\
\qquad  =\left(\sum\limits^{r_{(j+1)}}_{s=0}h_{i_{j+1},s}\bigl(q^{\frac{1}{2}}\bigr)\bigl(q^{\frac{1}{2d_{(j+1)}}}L_{j,j+1}\bigr)^{s}\right)^{-\varepsilon_{j+1}\{(\mathbf{g}_{k},d_{(j+1)}\mathbf{c}_{j+1}^{+})_{D}\}}\\
\qquad  =L_{j+1}^{-\varepsilon_{j+1}\{(\mathbf{g}_{k},d_{(j+1)}\mathbf{c}_{j+1}^{+})_{D}\}}.
 \end{gather*}
 It follows that
 \begin{align*}
 \mathbf{X}_{t_{0}}(-\mathbf{\tilde{g}}_{k})\mu_{t_{k}}^{t_{0}}(\mathbf{X}_{i_{k};t_{k}})
 ={}&\mathbf{X}_{t_{0}}(-\mathbf{\tilde{g}}_{k})\psi(\mathbf{i}_{1})(\mathbf{X}_{t_{0}}(\mathbf{\tilde{g}}_{k}))\psi(\mathbf{i}_{1})(\mathbf{X}_{t_{0}}(-\mathbf{\tilde{g}}_{k}))\psi(\mathbf{i}_{2})(\mathbf{X}_{t_{0}}(\mathbf{\tilde{g}}_{k}))\\
 &\cdots\psi(\mathbf{i}_{k-1})(\mathbf{X}_{t_{0}}(-\mathbf{\tilde{g}}_{k}))\psi(\mathbf{i}_{k})(\mathbf{X}_{t_{0}}(\mathbf{\tilde{g}}_{k}))\\
 ={}&L_{1}^{-\varepsilon_{1}\{(\mathbf{g}_{k},d_{(1)}\mathbf{c}_{1}^{+})_{D}\}}L_{2}^{-\varepsilon_{2}\{(\mathbf{g}_{k},d_{(2)}\mathbf{c}_{2}^{+})_{D}\}}\cdots L_{k}^{-\varepsilon_{k}\{(\mathbf{g}_{k},d_{(k)}\mathbf{c}_{k}^{+})_{D}\}}.\tag*{\qed}
 \end{align*}\renewcommand{\qed}{}
\end{proof}

\begin{Definition}
 There is a unique element $F_{i_k;t_k}:=F_{i_k;t_k}[\mathbf{Z}_1,\dots, \mathbf{Z}_n]$ in the skew field $\mathcal{F}_{DB}$ of~$\mathcal{T}_{DB}$ such that
 \begin{align}\label{eq:gupta-formula}
 \prod\limits^{\longrightarrow}_{j\in[1,k]}L_{j}^{-\varepsilon_{j}\{d_{(j)}(\mathbf{c}_{j}^{+},\mathbf{g}_{k})_D\}}=F_{i_k;t_k}\big[\hat{\mathbf{Y}}_{1;t_0},\dots, \hat{\mathbf{Y}}_{n;t_0}\big].
 \end{align}
 The element $F_{i_k,t_k}$ is called the {\it $F$-polynomial} of $\mathbf{X}_{t_k}(e_{i_k})$ whenever $F_{i_k;t_k}$ is a polynomial in $\mathbf{Z}_1,\dots,\mathbf{Z}_n$. With the help of $F_{i_k;t_k}$, we can rewrite \eqref{eq:f-poly-well-def} as
 \[
 \mathbf{X}_{i_k;t_k}=\mathbf{X}_{t_0}(\tilde{\mathbf{g}}_k)F_{i_k;t_k}\big[\hat{\mathbf{Y}}_{1;t_0},\dots, \hat{\mathbf{Y}}_{n;t_0}\big].
 \]

\end{Definition}
\begin{Remark}
 Recall that the $(\mathbf{r},\mathbf{z})$-$C$-pattern and $(\mathbf{r},\mathbf{z})$-$G$-pattern are uniquely determined by $B$, $R$ and $t_0\in \mathbb{T}_n$. It follows that the element $F_{i_k;t_k}$ only depends on $B$, $R$, $D$ and $t_0$.
\end{Remark}

The following is the main result of this section.
\begin{Theorem}\label{t:F-polynomial}\qquad
\begin{itemize}\itemsep=0pt
 \item[$(1)$] The element $F_{i_k;t_k}$ is a Laurent polynomial in $\mathbf{Z}_1,\dots, \mathbf{Z}_n$.
 \item[$(2)$] Suppose that $h_{i,s}(1)>0$ for each $i\in [1,n]$ and $s\in [1,r_i-1]$, then $F_{i_k;t_k}$ is a polynomial in $\mathbf{Z}_1,\dots, \mathbf{Z}_n$.
\end{itemize}

\end{Theorem}
\begin{proof}
 Since $F_{i_k;t_k}$ only depends on $B$, $D$ and $R$, but not on coefficients, it suffices to prove the statement for a particular choice of coefficients. Let \smash{$\tilde{B}^\bullet=\big[\begin{smallmatrix}
 B\\ I_n
 \end{smallmatrix}\big]$} be the $2n\times n$ matrix, where~$I_n$ is the identity matrix. There is a $(2n\times 2n)$-skew symmetric matrix $\Lambda^{\bullet}$ such that $\bigl(\tilde{B}^{\bullet},\Lambda^{\bullet}\bigr)$ is a compatible pair with \smash{$\bigl(\tilde{B}^{\bullet}\bigr)^{\mathsf T}\Lambda^{\bullet}=[D\quad 0]$} (cf.\ \cite[Example 0.5]{Berenstein_Zelevinsky2005}). It is clear that the \smash{$\mathbb{Z}\big[q^{\pm\frac{1}{2}}\big]$}-algebra generated by $\bigl\{\mathbf{\hat{Y}}_{t_0}^{\alpha}:=\mathbf{X}_{t_{0}}\bigl(\tilde{B}^{\bullet}\alpha\bigr)\mid \alpha\in\mathbb{Z}^{n}\bigr\}$ is isomorphic to $\mathcal{T}_{DB}$.

 According to Lemma \ref{product_of_F_polynomial}, there are two polynomials $A(\mathbf{Z}_{1},\dots,\mathbf{Z}_{n}),P(\mathbf{Z}_{1},\dots,\mathbf{Z}_{n})\in \mathcal{T}_{DB}$ with coefficients in \smash{$\mathbb{N}\big[q^{\pm\frac{1}{2}}\big]$} for \smash{$h_{i,s}\bigl(q^{\frac{1}{2}}\bigr)$}\footnote{Here we temporary regard $h_{i,s}\bigl(q^{\frac{1}{2}}\bigr)$ as a variable.} and $\mathbf{Z}_{1},\dots,\mathbf{Z}_{n}$ such that
 \begin{align}\label{fraction_of_F-polynomial}
 F_{i_{k};t_{k}}\bigl(\mathbf{\hat{Y}}_{1;t_{0}},\dots,\mathbf{\hat{Y}}_{n;t_{0}}\bigr)=A\bigl(\mathbf{\hat{Y}}_{1;t_{0}},\dots,\mathbf{\hat{Y}}_{n;t_{0}}\bigr)P\bigl(\mathbf{\hat{Y}}_{1;t_{0}},\dots,\mathbf{\hat{Y}}_{n;t_{0}}\bigr)^{-1}.
 \end{align}


 By \cite[Theorem 3.1]{BCDX2023}, $F_{i_{k};t_{k}}\bigl(\mathbf{\hat{Y}}_{1;t_{0}},\dots,\mathbf{\hat{Y}}_{n;t_{0}}\bigr)$ is also a Laurent polynomial in $\mathbf{X}_{1;t_{0}},\dots,\mathbf{X}_{2n;t_{0}}$. Hence all of \smash{$F_{i_{k};t_{k}}\bigl(\mathbf{\hat{Y}}_{1;t_{0}},\dots,\mathbf{\hat{Y}}_{n;t_{0}}\bigr)$}, \smash{$A\bigl(\mathbf{\hat{Y}}_{1;t_{0}},\dots,\mathbf{\hat{Y}}_{n;t_{0}}\bigr)$} and \smash{$P\bigl(\mathbf{\hat{Y}}_{1;t_{0}},\dots,\mathbf{\hat{Y}}_{n;t_{0}}\bigr)$} are Laurent polynomials in $\mathbf{X}_{1;t_{0}},\dots,\mathbf{X}_{2n;t_{0}}$. Taking the Newton polytopes of both sides of \eqref{fraction_of_F-polynomial} as the Laurent polynomials in $\mathbf{X}_{1;t_0},\dots, \mathbf{X}_{2n;t_0}$, we obtain
$\text{New}(F_{i_{k};t_{k}})+\text{New}(P)=\text{New}(A)$,
 where the sum is the Minkowski sum. By the definition of Minkowski sum, we conclude that every exponent vector of $F_{i_{k};t_{k}}$ is a $\mathbb{Q}$-linear combination of exponents vectors of $A$ and $P$.
 In particular, the exponent vectors of
 \smash{$F_{i_{k};t_{k}}\bigl(\mathbf{\hat{Y}}_{1;t_{0}},\dots,\mathbf{\hat{Y}}_{n;t_{0}}\bigr)$} as a Laurent polynomial in $\mathbf{X}_{1;t_0}$,$\dots$, $\mathbf{X}_{2n;t_0}$ can be expressed as $\mathbb{Q}$-linear combinations of $\tilde{B}^{\bullet}f_{1},\dots,\tilde{B}^{\bullet}f_{n}$. However, the last $n$ coordinates of the vectors $\tilde{B}^{\bullet}f_{i}$ is given by the standard basis $f_1,\dots, f_{n}\in\mathbb{Z}^{n}$, so each exponent vector must be a $\mathbb{Z}$-linear combination of $\tilde{B}^{\bullet}f_{1},\dots,\tilde{B}^{\bullet}f_{n}$. This prove that $F_{i_{k};t_{k}}\bigl(\mathbf{\hat{Y}}_{1;t_{0}},\dots,\mathbf{\hat{Y}}_{n;t_{0}}\bigr)$ is a Laurent polynomial in $\mathbf{\hat{Y}}_{1;t_{0}},\dots,\mathbf{\hat{Y}}_{n;t_{0}}$. Hence, $F_{i_k;t_k}[\mathbf{Z}_1,\dots, \mathbf{Z}_n]$ is a Laurent polynomial in $\mathbf{Z}_{1},\dots,\mathbf{Z}_{n}$.

 Since $A(\mathbf{Z}_{1},\dots,\mathbf{Z}_{n}),P(\mathbf{Z}_{1},\dots,\mathbf{Z}_{n})$ have coefficients in $\mathbb{N}\big[q^{\pm\frac{1}{2}}\big]$ for $h_{i,s}\bigl(q^{\frac{1}{2}}\bigr)$ and $\mathbf{Z}_{1},\dots,\mathbf{Z}_{n}$. By the assumption $h_{i,s}(1)>0$ for $i\in[1,n]$ and $s\in[1,r_{i}-1]$, by setting \smash{$q^{\frac{1}{2}}=1$} does not shrink $\text{New}\bigl(F_{i_{k};t_{k}}\bigl(\mathbf{\hat{Y}}_{1;t_{0}},\dots,\mathbf{\hat{Y}}_{n;t_{0}}\bigr)\bigr)$. That is
 \[
 \text{New}\bigl(F_{i_{k};t_{k}}\bigl(\mathbf{\hat{Y}}_{1;t_{0}},\dots,\mathbf{\hat{Y}}_{n;t_{0}}\bigr)\bigr)
 =\text{New}\bigl(F_{i_{k};t_{k}}\bigl(\mathbf{\hat{Y}}_{1;t_{0}},\dots,\mathbf{\hat{Y}}_{n;t_{0}}\bigr)|_{q^{\frac{1}{2}}=1}\bigr).\]
By Proposition \ref{p:separation-formula-gca} and Lemma \ref{l:g-vector-comparison}, we conclude that \[F_{i_{k};t_{k}}\bigl(\mathbf{\hat{Y}}_{1;t_{0}},\dots,\mathbf{\hat{Y}}_{n;t_{0}}\bigr)|_{q^{\frac{1}{2}}=1}=F_{i_k;t_k}(\mathbf{\hat{y}},\mathbf{z})|_{z_{i,s}=h_{i,s}(1),i\in [1,n],s\in [1,r_i-1]},\] where $F_{i_k;t_k}(\mathbf{y},\mathbf{z})$ is the $F$-polynomial of the cluster variable $x_{i_k;t_k}$ of the corresponding generalized cluster algebra $\mathcal{A}^\bullet$ with principal coefficients.
 Thus $\text{New}\bigl(F_{i_{k};t_{k}}\bigl(\mathbf{\hat{Y}}_{1;t_{0}},\dots,\mathbf{\hat{Y}}_{n;t_{0}}\bigr)\bigr)$ does not contain any points with negative coordinates. It follows that \smash{$F_{i_{k};t_{k}}\bigl(\mathbf{\hat{Y}}_{1;t_{0}},\dots,\mathbf{\hat{Y}}_{n;t_{0}}\bigr)$} is a~polynomial in $\mathbf{\hat{Y}}_{1;t_{0}},\dots,\mathbf{\hat{Y}}_{n;t_{0}}$. Hence, $F_{i_{k};t_{k}}[\mathbf{Z}_1,\dots,\mathbf{Z}_{n}]$ is a polynomial in $\mathbf{Z}_{1},\dots,\mathbf{Z}_{n}$.
\end{proof}
